
\Data\ID{1003055601}
\Updated{2022-04-23 05:04:02}
\Level{A}
\SubArea{Logic and Sets}
\LANGen{
\Question{
Given the sets \( A=[ -12;12 ] \) and \( B=(3;20) \) find the set difference \( A\setminus B \).}
{\( [ -12;3 ] \)}{\( ( -12;3 ] \)}{\( [ -12;3 ) \)}{\( ( 12;20 ) \)}{}{}}
\LANGpl{
\Question{
Różnicą zbiorów \( A\setminus B \) dla  \( A=\langle -12;12 \rangle \), \( B=(3;20) \) jest:}
{\( \langle -12;3 \rangle \)}{\( ( -12;3 \rangle \)}{\( \langle -12;3 ) \)}{\( ( 12;20 ) \)}{}{}}
\LANGcs{
\Question{
Určete rozdíl množin \( A\setminus B \), jestliže  \( A=\langle -12;12 \rangle \), \( B=(3;20) \).}
{\( \langle -12;3 \rangle \)}{\( ( -12;3 \rangle \)}{\( \langle -12;3 ) \)}{\( ( 12;20 ) \)}{}{}}
\LANGsk{
\Question{
Určte rozdiel množín \( A\setminus B \), ak  \( A=\langle -12;12 \rangle \), \( B=(3;20) \).}
{\( \langle -12;3 \rangle \)}{\( ( -12;3 \rangle \)}{\( \langle -12;3 ) \)}{\( ( 12;20 ) \)}{}{}}
\LANGes{
\Question{
Calcula la diferencia de conjuntos  \( A\setminus B \) para  \( A=[ -12;12 ] \), \( B=(3;20) \).}
{\( [ -12;3 ] \)}{\( ( -12;3 ] \)}{\( [ -12;3 ) \)}{\( ( 12;20 ) \)}{}{}}
\End

\Data\ID{1003055602}
\Updated{2020-07-15 03:07:12}
\Level{A}
\SubArea{Logic and Sets}
\LANGen{
\Question{
Give the set difference \( A\setminus B \) for  \( A = [ -5;7 ] \), \( B=\{7;11\} \).}
{\( [ -5;7 ) \)}{\( [ -5;11) \)}{\( [ -5;7 )\cup(7;11) \)}{\( [ -5;7 )\cup\{11\} \)}{}{}}
\LANGpl{
\Question{
Różnicą zbiorów  \( A\setminus B \) dla  \( A = \langle -5;7 \rangle \), \( B=\{7;11\} \) jest:}
{\( \langle -5;7 ) \)}{\( \langle -5;11) \)}{\( \langle -5;7 )\cup(7;11) \)}{\( \langle -5;7 )\cup\{11\} \)}{}{}}
\LANGcs{
\Question{
Určete rozdíl množin \( A\setminus B \), jestliže  \( A = \langle -5;7 \rangle \), \( B=\{7;11\} \).}
{\( \langle -5;7 ) \)}{\( \langle -5;11) \)}{\( \langle -5;7 )\cup(7;11) \)}{\( \langle -5;7 )\cup\{11\} \)}{}{}}
\LANGsk{
\Question{
Určte rozdiel množín \( A\setminus B \), ak  \( A = \langle -5;7 \rangle \), \( B=\{7;11\} \).}
{\( \langle -5;7 ) \)}{\( \langle -5;11) \)}{\( \langle -5;7 )\cup(7;11) \)}{\( \langle -5;7 )\cup\{11\} \)}{}{}}
\LANGes{
\Question{
Calcula la diferencia de conjuntos  \( A\setminus B \) para  \( A = [ -5;7 ] \), \( B=\{7;11\} \).}
{\( [ -5;7 ) \)}{\( [ -5;11) \)}{\( [ -5;7 )\cup(7;11) \)}{\( [ -5;7 )\cup\{11\} \)}{}{}}
\End

\Data\ID{1003055603}
\Updated{2022-04-23 05:04:08}
\Level{A}
\SubArea{Logic and Sets}
\LANGen{
\Question{
Find the set difference \( A\setminus B \) for  \( A=\left\{x\in \mathbb{Z}\colon x^2=4\right\} \) and \( B=\{-1;0;1;2;3\} \).}
{\( \{-2\} \)}{\( \{-1;0;1;3\} \)}{\( \{-2;2\} \)}{\( \{-2;-1;0;1;2;3\} \)}{}{}}
\LANGpl{
\Question{
Różnicą zbiorów \( A\setminus B \) dla  \( A=\left\{x\in \mathbb{Z}\colon x^2=4\right\} \), \( B=\{-1;0;1;2;3\} \) jest:}
{\( \{-2\} \)}{\( \{-1;0;1;3\} \)}{\( \{-2;2\} \)}{\( \{-2;-1;0;1;2;3\} \)}{}{}}
\LANGcs{
\Question{
Určete rozdíl množin \( A\setminus B \) jestliže  \( A=\left\{x\in \mathbb{Z}\colon x^2=4\right\} \), \( B=\{-1;0;1;2;3\} \).}
{\( \{-2\} \)}{\( \{-1;0;1;3\} \)}{\( \{-2;2\} \)}{\( \{-2;-1;0;1;2;3\} \)}{}{}}
\LANGsk{
\Question{
Určte rozdiel množín \( A\setminus B \) ak  \( A=\left\{x\in \mathbb{Z}\colon x^2=4\right\} \), \( B=\{-1;0;1;2;3\} \).}
{\( \{-2\} \)}{\( \{-1;0;1;3\} \)}{\( \{-2;2\} \)}{\( \{-2;-1;0;1;2;3\} \)}{}{}}
\LANGes{
\Question{
Calcula la diferencia de conjuntos \( A\setminus B \) para  \( A=\left\{x\in \mathbb{Z}\colon x^2=4\right\} \), \( B=\{-1;0;1;2;3\} \).}
{\( \{-2\} \)}{\( \{-1;0;1;3\} \)}{\( \{-2;2\} \)}{\( \{-2;-1;0;1;2;3\} \)}{}{}}
\End

\Data\ID{1003055604}
\Updated{2020-07-15 03:07:12}
\Level{A}
\SubArea{Logic and Sets}
\LANGen{
\Question{
Let  \( A=\{0;1;2;3\} \), \( B=\{0;1;2\} \) and \( C=\{x\colon x=3k+l\} \), where \( k\in A\), \(l\in B\). Which of the following sets specifies \( C \) by the list of all its elements?}
{\( \{0;1;2;3;4;5;6;7;8;9;10;11\} \)}{\( \{4;5;6;7;8;9;10;11\} \)}{\( [ 0;11 ] \)}{\( \{0;1;2;3\} \)}{}{}}
\LANGpl{
\Question{
Niech  \( A=\{0;1;2;3\} \) i \( B=\{0;1;2\} \). Zbiór  \( C=\{x\colon x=3k+l\} \), gdzie \( k\in A\), \(l\in B\).  Który z poniższych zbiorów określa zbiór elementów zbioru  \( C \)?}
{\( \{0;1;2;3;4;5;6;7;8;9;10;11\} \)}{\( \{4;5;6;7;8;9;10;11\} \)}{\( \langle 0;11 \rangle \)}{\( \{0;1;2;3\} \)}{}{}}
\LANGcs{
\Question{
Nechť  \( A=\{0;1;2;3\} \) a \( B=\{0;1;2\} \). Dále nechť  \( C=\{x \in\mathbb{Z}\colon x=3k+l, k\in A, l\in B \} \). Určete, která z následujících množin daných výčtem prvků, je rovna množině \( C \) ?}
{\( \{0;1;2;3;4;5;6;7;8;9;10;11\} \)}{\( \{4;5;6;7;8;9;10;11\} \)}{\( \langle 0;11 \rangle \)}{\( \{0;1;2;3\} \)}{}{}}
\LANGsk{
\Question{
Nech  \( A=\{0;1;2;3\} \) a \( B=\{0;1;2\} \). Ak množina  \( C=\{x\colon x=3k+l\} \), kde \( k\in A\), \(l\in B\) ktorá z následujúcich množin určuje \( C \) množinu všetkých svojich prvkov?}
{\( \{0;1;2;3;4;5;6;7;8;9;10;11\} \)}{\( \{4;5;6;7;8;9;10;11\} \)}{\( \langle 0;11 \rangle \)}{\( \{0;1;2;3\} \)}{}{}}
\LANGes{
\Question{
Sean  \( A=\{0;1;2;3\} \), \( B=\{0;1;2\} \) y \( C=\{x\colon x=3k+l\} \), donde \( k\in A\), \(l\in B\). ¿Cuál de los siguientes conjuntos define \( C \) dando la lista de todos sus elementos?}
{\( \{0;1;2;3;4;5;6;7;8;9;10;11\} \)}{\( \{4;5;6;7;8;9;10;11\} \)}{\( [ 0;11 ] \)}{\( \{0;1;2;3\} \)}{}{}}
\End

\Data\ID{1003055610}
\Updated{2020-07-15 03:07:12}
\Level{A}
\SubArea{Logic and Sets}
\LANGen{
\Question{
Let \( A=(-3;5] \) and \( B=[-1;+\infty) \). Find the intersection \( A\cap B \).}
{\( [ -1;5] \)}{\( (-1;5) \)}{\( (-3;+\infty) \)}{\( (-1;5] \)}{}{}}
\LANGpl{
\Question{
Jeżeli \( A=(-3;5\rangle \)  i  \( B=\langle-1;+\infty) \). To \( A\cap B \) jest równy:}
{\( \langle -1;5\rangle \)}{\( (-1;5) \)}{\( (-3;+\infty) \)}{\( (-1;5\rangle \)}{}{}}
\LANGcs{
\Question{
Nechť \( A=(-3;5\rangle \) a \( B=\langle-1;+\infty) \). Určete průnik \( A\cap B \).}
{\( \langle -1;5\rangle \)}{\( (-1;5) \)}{\( (-3;+\infty) \)}{\( (-1;5\rangle \)}{}{}}
\LANGsk{
\Question{
Nech \( A=(-3;5\rangle \) a \( B=\langle-1;+\infty) \). Určte prienik \( A\cap B \).}
{\( \langle -1;5\rangle \)}{\( (-1;5) \)}{\( (-3;+\infty) \)}{\( (-1;5\rangle \)}{}{}}
\LANGes{
\Question{
Sean \( A=(-3;5] \) y \( B=[-1;+\infty) \). Calcula la intersección  \( A\cap B \).}
{\( [ -1;5] \)}{\( (-1;5) \)}{\( (-3;+\infty) \)}{\( (-1;5] \)}{}{}}
\End

\Data\ID{1003055611}
\Updated{2020-07-15 03:07:12}
\Level{A}
\SubArea{Logic and Sets}
\LANGen{
\Question{
Let \( A=(-1;5] \) and \( B=[ -1;7) \). Give the union \( A\cup B\).}
{\( [ -1;7 ) \)}{\( ( -1;7 ) \)}{\( ( -1;5 ) \)}{\( [ -1;5] \)}{}{}}
\LANGpl{
\Question{
Jeżeli \( A=(-1;5\rangle \) i \( B=\langle -1;7) \). To \( A\cup B\) jest równe:}
{\( \langle -1;7 ) \)}{\( ( -1;7 ) \)}{\( ( -1;5 ) \)}{\( \langle -1;5\rangle \)}{}{}}
\LANGcs{
\Question{
Nechť \( A=(-1;5\rangle \) a \( B=\langle -1;7) \). Urči sjednocení \( A\cup B\).}
{\( \langle -1;7 ) \)}{\( ( -1;7 ) \)}{\( ( -1;5 ) \)}{\( \langle -1;5\rangle \)}{}{}}
\LANGsk{
\Question{
Nech \( A=(-1;5\rangle \) a \( B=\langle -1;7) \). Urč zjednotenie \( A\cup B\).}
{\( \langle -1;7 ) \)}{\( ( -1;7 ) \)}{\( ( -1;5 ) \)}{\( \langle -1;5\rangle \)}{}{}}
\LANGes{
\Question{
Sean \( A=(-1;5] \) y \( B=[ -1;7) \). Calcula la unión  \( A\cup B\).}
{\( [ -1;7 ) \)}{\( ( -1;7 ) \)}{\( ( -1;5 ) \)}{\( [ -1;5] \)}{}{}}
\End

\Data\ID{1003055612}
\Updated{2022-04-23 05:04:08}
\Level{A}
\SubArea{Logic and Sets}
\LANGen{
\Question{
Find the intersection \( A\cap B' \) if \( A=(-4;+\infty) \) and \(B=(-\infty;6) \). (By \(B'\) the complement of the set \( B \) is denoted.)}
{\( [ 6;+\infty) \)}{\( (-4;6) \)}{\( [-4;6] \)}{\( (-\infty;4] \)}{}{}}
\LANGpl{
\Question{
Wyznacz \( A\cap B' \) jeżeli \( A=(-4;+\infty) \) i \(B=(-\infty;6) \). (Przez \(B'\) oznaczamy dopełnienie zbioru \( B \).)}
{\( \langle 6;+\infty) \)}{\( (-4;6) \)}{\( \langle-4;6\rangle \)}{\( (-\infty;4\rangle \)}{}{}}
\LANGcs{
\Question{
Najdi průnik množiny \( A\cap B' \), jestliže \( A=(-4;+\infty) \) a \(B=(-\infty;6) \). (\(B'\) označuje doplněk množiny \( B \).)}
{\( \langle 6;+\infty) \)}{\( (-4;6) \)}{\( \langle-4;6\rangle \)}{\( (-\infty;4\rangle \)}{}{}}
\LANGsk{
\Question{
Nájdi prienik množiny \( A\cap B' \), ak \( A=(-4;+\infty) \) a \(B=(-\infty;6) \). (\(B'\) označuje doplnok množiny \( B \).)}
{\( \langle 6;+\infty) \)}{\( (-4;6) \)}{\( \langle-4;6\rangle \)}{\( (-\infty;4\rangle \)}{}{}}
\LANGes{
\Question{
Calcula la intersección  \( A\cap B' \) si \( A=(-4;+\infty) \) y \(B=(-\infty;6) \). (Por \(B'\) se denota el complementario del conjunto \( B \).)}
{\( [ 6;+\infty) \)}{\( (-4;6) \)}{\( [-4;6] \)}{\( (-\infty;4] \)}{}{}}
\End

\Data\ID{1003055613}
\Updated{2020-07-15 03:07:12}
\Level{A}
\SubArea{Logic and Sets}
\LANGen{
\Question{
Give the intersection \( A\cap B \) if \( A=[-7;1] \) and \( B=(1;2) \).}
{\( \emptyset \)}{\( \{1\} \)}{\( [-7;2) \)}{\( (-7;2) \)}{}{}}
\LANGpl{
\Question{
Wyznacz \( A\cap B \),  jeśli \( A=\langle-7;1\rangle \) i \( B=(1;2) \).}
{\( \emptyset \)}{\( \{1\} \)}{\( \langle-7;2) \)}{\( (-7;2) \)}{}{}}
\LANGcs{
\Question{
Určete průnik \( A\cap B \), jestliže \( A=\langle-7;1\rangle \) a \( B=(1;2) \).}
{\( \emptyset \)}{\( \{1\} \)}{\( \langle-7;2) \)}{\( (-7;2) \)}{}{}}
\LANGsk{
\Question{
Určte prienik \( A\cap B \) ak \( A=\langle-7;1\rangle \) a \( B=(1;2) \).}
{\( \emptyset \)}{\( \{1\} \)}{\( \langle-7;2) \)}{\( (-7;2) \)}{}{}}
\LANGes{
\Question{
Calcula la intersección  \( A\cap B \) si \( A=[-7;1] \) y \( B=(1;2) \).}
{\( \emptyset \)}{\( \{1\} \)}{\( [-7;2) \)}{\( (-7;2) \)}{}{}}
\End

\Data\ID{1003055702}
\Updated{2020-07-15 03:07:12}
\Level{A}
\SubArea{Logic and Sets}
\LANGen{
\Question{
The set of all numbers that satisfy the  following relations
\[ (x \geq -1) \wedge (x > -2) \wedge (x < 3) \] 
can be written as:}
{\( [ -1;3 ) \)}{\( \mathbb{R} \)}{\( [ -2;3) \)}{\( (-2;-1] \)}{}{}}
\LANGpl{
\Question{
Liczby spełniające układ: 
\[ (x \geq -1) \wedge (x > -2) \wedge (x < 3) \] 
można zapisać}
{\( \langle -1;3 ) \)}{\( \mathbb{R} \)}{\( \langle -2;3) \)}{\( (-2;-1\rangle \)}{}{}}
\LANGcs{
\Question{
Množina všech čísel, která splňují následující vztahy,
\[ (x \geq -1) \wedge (x > -2) \wedge (x < 3) \] 
může být zapsána jako:}
{\( \langle -1;3 ) \)}{\( \mathbb{R} \)}{\( \langle -2;3) \)}{\( (-2;-1\rangle \)}{}{}}
\LANGsk{
\Question{
Množina všetkých čísel, ktoré spĺňajú nasledujúce vzťahy,
\[ (x \geq -1) \wedge (x > -2) \wedge (x < 3) \] 
môže byť zapísaná ako}
{\( \langle -1;3 ) \)}{\( \mathbb{R} \)}{\( \langle -2;3) \)}{\( (-2;-1\rangle \)}{}{}}
\LANGes{
\Question{
El conjunto de todos los números que satisfacen las siguientes relaciones 
\[ (x \geq -1) \wedge (x > -2) \wedge (x < 3) \] 
se puede escribir como:}
{\( [ -1;3 ) \)}{\( \mathbb{R} \)}{\( [ -2;3) \)}{\( (-2;-1] \)}{}{}}
\End

\Data\ID{1003055704}
\Updated{2020-07-15 03:07:12}
\Level{A}
\SubArea{Logic and Sets}
\LANGen{
\Question{
Find the interval that is the difference of sets \( A \) and \( B \) if \( A = [ -8; 12] \) and \( B = (0; 20) \).}
{\( [-8;0] \)}{\( (-8;0) \)}{\( [-8;0) \)}{\( (-8;0] \)}{}{}}
\LANGpl{
\Question{
Wyznacz różnicę zbiorów \( A = \langle-8; 12\rangle \) i \( B = (0; 20) \).}
{\( \langle-8;0\rangle \)}{\( (-8;0) \)}{\( \langle-8;0) \)}{\( (-8;0\rangle \)}{}{}}
\LANGcs{
\Question{
Urči interval, který je rozdílem množin \( A \) a \( B \), jestliže \( A = \langle -8; 12\rangle \) a \( B = (0; 20) \).}
{\( \langle-8;0\rangle \)}{\( (-8;0) \)}{\( \langle-8;0) \)}{\( (-8;0\rangle \)}{}{}}
\LANGsk{
\Question{
Urč interval, ktorý je rozdielom množín \( A \) a \( B \), ak \( A = \langle -8; 12\rangle \) a \( B = (0; 20) \).}
{\( \langle-8;0\rangle \)}{\( (-8;0) \)}{\( \langle-8;0) \)}{\( (-8;0\rangle \)}{}{}}
\LANGes{
\Question{
Calcula el intervalo que es la diferencia de los conjuntos  \( A \) y \( B \) si \( A = [ -8; 12] \) y \( B = (0; 20) \).}
{\( [-8;0] \)}{\( (-8;0) \)}{\( [-8;0) \)}{\( (-8;0] \)}{}{}}
\End

\Data\ID{1003055705}
\Updated{2022-04-23 05:04:57}
\Level{A}
\SubArea{Logic and Sets}
\LANGen{
\Question{
Given the  sets
\begin{gather*} A = \{1; 2; 3; 4; 5; 6; 7; 8; 9\},\\ B = \{3; 4; 5; 6; 7\},\\ C = \{6; 7; 8; 9; 10; 11; 12\}, \end{gather*}
find the union  \( A\cup B\cup C \).}
{\( \{1; 2; 3; 4; 5; 6; 7; 8; 9;10; 11; 12\} \)}{\( \{1; 2; 3; 4; 5; 6; 7; 8; 9\} \)}{\( \{6; 7\} \)}{\( \{1; 2; 10; 11; 12\} \)}{}{}}
\LANGpl{
\Question{
Dane są zbiory:
\begin{gather*} A = \{1; 2; 3; 4; 5; 6; 7; 8; 9\},\\ B = \{3; 4; 5; 6; 7\},\\ C = \{6; 7; 8; 9; 10; 11; 12\}, \end{gather*}
sumą zbiorów  \( A\cup B\cup C \) jest:}
{\( \{1; 2; 3; 4; 5; 6; 7; 8; 9;10; 11; 12\} \)}{\( \{1; 2; 3; 4; 5; 6; 7; 8; 9\} \)}{\( \{6; 7\} \)}{\( \{1; 2; 10; 11; 12\} \)}{}{}}
\LANGcs{
\Question{
Máme dány množiny
\begin{gather*} A = \{1; 2; 3; 4; 5; 6; 7; 8; 9\},\\ B = \{3; 4; 5; 6; 7\},\\ C = \{6; 7; 8; 9; 10; 11; 12\}, \end{gather*}
najdi sjednocení  \( A\cup B\cup C \).}
{\( \{1; 2; 3; 4; 5; 6; 7; 8; 9;10; 11; 12\} \)}{\( \{1; 2; 3; 4; 5; 6; 7; 8; 9\} \)}{\( \{6; 7\} \)}{\( \{1; 2; 10; 11; 12\} \)}{}{}}
\LANGsk{
\Question{
Máme dané množiny
\begin{gather*} A = \{1; 2; 3; 4; 5; 6; 7; 8; 9\},\\ B = \{3; 4; 5; 6; 7\},\\ C = \{6; 7; 8; 9; 10; 11; 12\}, \end{gather*}
nájdi zjednotenie  \( A\cup B\cup C \).}
{\( \{1; 2; 3; 4; 5; 6; 7; 8; 9;10; 11; 12\} \)}{\( \{1; 2; 3; 4; 5; 6; 7; 8; 9\} \)}{\( \{6; 7\} \)}{\( \{1; 2; 10; 11; 12\} \)}{}{}}
\LANGes{
\Question{
Dados los conjuntos 
\begin{gather*} A = \{1; 2; 3; 4; 5; 6; 7; 8; 9\},\\ B = \{3; 4; 5; 6; 7\},\\ C = \{6; 7; 8; 9; 10; 11; 12\}, \end{gather*}
calcula la unión \( A\cup B\cup C \).}
{\( \{1; 2; 3; 4; 5; 6; 7; 8; 9;10; 11; 12\} \)}{\( \{1; 2; 3; 4; 5; 6; 7; 8; 9\} \)}{\( \{6; 7\} \)}{\( \{1; 2; 10; 11; 12\} \)}{}{}}
\End

\Data\ID{1003055706}
\Updated{2021-01-05 12:01:51}
\Level{A}
\SubArea{Logic and Sets}
\LANGen{
\Question{
Given the  sets 
\begin{gather*} A = \{ 3; 4; 5; 6; 7; 8\},\\ B = \{3; 4; 5; 6; 7\},\\ C = \{6; 7; 8; 9; 10;11\}, \end{gather*}
find the intersection \( A\cap B\cap C \).}
{\( \{ 6; 7\} \)}{\( \{3; 4; 5; 6; 7; 8; 9; 10; 11\} \)}{\( \{3; 4; 5; 6; 7\} \)}{\( \{9; 10; 11\} \)}{}{}}
\LANGpl{
\Question{
Dane są zbiory: 
\begin{gather*} A = \{ 3; 4; 5; 6; 7; 8\},\\ B = \{3; 4; 5; 6; 7\},\\ C = \{6; 7; 8; 9; 10;11\}, \end{gather*}
iloczynem zbiorów \( A\cap B\cap C \) jest:}
{\( \{ 6; 7\} \)}{\( \{3; 4; 5; 6; 7; 8; 9; 10; 11\} \)}{\( \{3; 4; 5; 6; 7\} \)}{\( \{9; 10; 11\} \)}{}{}}
\LANGcs{
\Question{
Máme dány množiny 
\begin{gather*} A = \{ 3; 4; 5; 6; 7; 8\},\\ B = \{3; 4; 5; 6; 7\},\\ C = \{6; 7; 8; 9; 10;11\}, \end{gather*}
najdi průnik \( A\cap B\cap C \).}
{\( \{ 6; 7\} \)}{\( \{3; 4; 5; 6; 7; 8; 9; 10; 11\} \)}{\( \{3; 4; 5; 6; 7\} \)}{\( \{9; 10; 11\} \)}{}{}}
\LANGsk{
\Question{
Máme dané množiny 
\begin{gather*} A = \{ 3; 4; 5; 6; 7; 8\},\\ B = \{3; 4; 5; 6; 7\},\\ C = \{6; 7; 8; 9; 10;11\}, \end{gather*}
nájdi prienik \( A\cap B\cap C \).}
{\( \{ 6; 7\} \)}{\( \{3; 4; 5; 6; 7; 8; 9; 10; 11\} \)}{\( \{3; 4; 5; 6; 7\} \)}{\( \{9; 10; 11\} \)}{}{}}
\LANGes{
\Question{
Dados los conjuntos 
\begin{gather*} A = \{ 3; 4; 5; 6; 7; 8\},\\ B = \{3; 4; 5; 6; 7\},\\ C = \{6; 7; 8; 9; 10;11\}, \end{gather*}
calcula la intersección \( A\cap B\cap C \).}
{\( \{ 6; 7\} \)}{\( \{3; 4; 5; 6; 7; 8; 9; 10; 11\} \)}{\( \{3; 4; 5; 6; 7\} \)}{\( \{9; 10; 11\} \)}{}{}}
\End

\Data\ID{1003055707}
\Updated{2020-07-15 03:07:12}
\Level{A}
\SubArea{Logic and Sets}
\LANGen{
\Question{
Let \( A = \{a; b; c; d; e; f\} \), \( B = \{a; k; b; l\} \). Find the union \( A\cup B \).}
{\( \{a; b; c; d; e; f; k; l\} \)}{\( \{a; b\} \)}{\( \{ c; d; e; f\} \)}{\( \{k; l\} \)}{}{}}
\LANGpl{
\Question{
Dane są zbiory: \( A = \{a; b; c; d; e; f\} \), \( B = \{a; k; b; l\} \). Sumą zbiorów \( A\cup B \).}
{\( \{a; b; c; d; e; f; k; l\} \)}{\( \{a; b\} \)}{\( \{ c; d; e; f\} \)}{\( \{k; l\} \)}{}{}}
\LANGcs{
\Question{
Nechť \( A = \{a; b; c; d; e; f\} \), \( B = \{a; k; b; l\} \). Najdi sjednocení \( A\cup B \).}
{\( \{a; b; c; d; e; f; k; l\} \)}{\( \{a; b\} \)}{\( \{ c; d; e; f\} \)}{\( \{k; l\} \)}{}{}}
\LANGsk{
\Question{
Nech \( A = \{a; b; c; d; e; f\} \), \( B = \{a; k; b; l\} \). Nájdi zjednotenie \( A\cup B \).}
{\( \{a; b; c; d; e; f; k; l\} \)}{\( \{a; b\} \)}{\( \{ c; d; e; f\} \)}{\( \{k; l\} \)}{}{}}
\LANGes{
\Question{
Sean \( A = \{a; b; c; d; e; f\} \), \( B = \{a; k; b; l\} \). Calcula la unión \( A\cup B \).}
{\( \{a; b; c; d; e; f; k; l\} \)}{\( \{a; b\} \)}{\( \{ c; d; e; f\} \)}{\( \{k; l\} \)}{}{}}
\End

\Data\ID{1003055708}
\Updated{2021-07-24 10:07:20}
\Level{A}
\SubArea{Logic and Sets}
\LANGen{
\Question{
Let \( A = \{a; b; c; d; e; f\} \), \( B = \{a; g; b; h; i\} \). Determine the intersection \( A\cap B\).}
{\( \{a; b\} \)}{\( \{a; b; c; d; e; f; g; h; i\} \)}{\( \{ g; h; i\} \)}{\( \{c; d; e; f\} \)}{}{}}
\LANGpl{
\Question{
Dane są zbiory: \( A = \{a; b; c; d; e; f\} \), \( B = \{a; g; b; h; i\} \). Iloczynem zbiorów \( A\cap B\) jest:}
{\( \{a; b\} \)}{\( \{a; b; c; d; e; f; g; h; i\} \)}{\( \{ g; h; i\} \)}{\( \{c; d; e; f\} \)}{}{}}
\LANGcs{
\Question{
Nechť \( A = \{a; b; c; d; e; f\} \), \( B = \{a; g; b; h; i\} \). Urči průnik \( A\cap B\).}
{\( \{a; b\} \)}{\( \{a; b; c; d; e; f; g; h; i\} \)}{\( \{ g; h; i\} \)}{\( \{c; d; e; f\} \)}{}{}}
\LANGsk{
\Question{
Nech \( A = \{a; b; c; d; e; f\} \), \( B = \{a; g; b; h; i\} \). Urč prienik \( A\cap B\).}
{\( \{a; b\} \)}{\( \{a; b; c; d; e; f; g; h; i\} \)}{\( \{ g; h; i\} \)}{\( \{c; d; e; f\} \)}{}{}}
\LANGes{
\Question{
Sean \( A = \{a; b; c; d; e; f\} \), \( B = \{a; g; b; h; i\} \). Calcula la intersección \( A\cap B\).}
{\( \{a; b\} \)}{\( \{a; b; c; d; e; f; g; h; i\} \)}{\( \{ g; h; i\} \)}{\( \{c; d; e; f\} \)}{}{}}
\End

\Data\ID{1003055709}
\Updated{2020-07-15 03:07:12}
\Level{A}
\SubArea{Logic and Sets}
\LANGen{
\Question{
Give the intersection \( A\cap B \) if \( A=\{1; 3; 5; 7; 9\} \) and  \( B=\{0; 3; 6; 9\} \).}
{\( \{3; 9\} \)}{\( \{0; 1; 3; 5; 6; 7; 9\} \)}{\( \{1; 5; 7\} \)}{\( \{0; 6\} \)}{}{}}
\LANGpl{
\Question{
Iloczynem zbiorów \( A\cap B \) dla \( A=\{1; 3; 5; 7; 9\} \) i  \( B=\{0; 3; 6; 9\} \) jest zbiór:}
{\( \{3; 9\} \)}{\( \{0; 1; 3; 5; 6; 7; 9\} \)}{\( \{1; 5; 7\} \)}{\( \{0; 6\} \)}{}{}}
\LANGcs{
\Question{
Urči průnik \( A\cap B \), jestliže \( A=\{1; 3; 5; 7; 9\} \) a  \( B=\{0; 3; 6; 9\} \).}
{\( \{3; 9\} \)}{\( \{0; 1; 3; 5; 6; 7; 9\} \)}{\( \{1; 5; 7\} \)}{\( \{0; 6\} \)}{}{}}
\LANGsk{
\Question{
Urč prienik \( A\cap B \), ak \( A=\{1; 3; 5; 7; 9\} \) a  \( B=\{0; 3; 6; 9\} \).}
{\( \{3; 9\} \)}{\( \{0; 1; 3; 5; 6; 7; 9\} \)}{\( \{1; 5; 7\} \)}{\( \{0; 6\} \)}{}{}}
\LANGes{
\Question{
Calcula la intersección \( A\cap B \) si \( A=\{1; 3; 5; 7; 9\} \) y  \( B=\{0; 3; 6; 9\} \).}
{\( \{3; 9\} \)}{\( \{0; 1; 3; 5; 6; 7; 9\} \)}{\( \{1; 5; 7\} \)}{\( \{0; 6\} \)}{}{}}
\End

\Data\ID{1003055710}
\Updated{2020-07-15 03:07:12}
\Level{A}
\SubArea{Logic and Sets}
\LANGen{
\Question{
Give the difference of sets \( A \) and \( B \) if \( A = \{4; 5; 6\}\) and \(B = \{5; 6; 7\} \).}
{\( \{4\} \)}{\( \{7\} \)}{\( \{4; 5; 6; 7\} \)}{\( [ 4;5 ] \)}{}{}}
\LANGpl{
\Question{
Wyznacz różnicę zbiorów \( A \) i \( B \), jeśli \( A = \{4; 5; 6\}\) i \(B = \{5; 6; 7\} \).}
{\( \{4\} \)}{\( \{7\} \)}{\( \{4; 5; 6; 7\} \)}{\( \langle 4;5 \rangle \)}{}{}}
\LANGcs{
\Question{
Urči rozdíl množin \( A \) a \( B \), kde \( A = \{4; 5; 6\}\) a \(B = \{5; 6; 7\} \).}
{\( \{4\} \)}{\( \{7\} \)}{\( \{4; 5; 6; 7\} \)}{\( \langle 4;5 \rangle \)}{}{}}
\LANGsk{
\Question{
Urč rozdiel množín \( A \) a \( B \), kde \( A = \{4; 5; 6\}\) a \(B = \{5; 6; 7\} \).}
{\( \{4\} \)}{\( \{7\} \)}{\( \{4; 5; 6; 7\} \)}{\( \langle 4;5 \rangle \)}{}{}}
\LANGes{
\Question{
Calcula la diferencia de los conjuntos \( A \) y \( B \) si \( A = \{4; 5; 6\}\) y \(B = \{5; 6; 7\} \).}
{\( \{4\} \)}{\( \{7\} \)}{\( \{4; 5; 6; 7\} \)}{\( [ 4;5 ] \)}{}{}}
\End

\Data\ID{1003055711}
\Updated{2020-07-15 03:07:12}
\Level{A}
\SubArea{Logic and Sets}
\LANGen{
\Question{
Determine the set  of all common divisors of numbers \( 24 \) and \( 18 \).}
{\( \{ 1;2;3;6 \} \)}{\( \{ 1;2;3;4;6;8;9;12;18;24 \} \)}{\( \{ 2;3;6 \} \)}{\( \{ 1;2;3;4;6;9;12 \} \)}{}{}}
\LANGpl{
\Question{
Zbiór wspólnych dzielników liczb  \( 24 \) i \( 18 \).}
{\( \{ 1;2;3;6 \} \)}{\( \{ 1;2;3;4;6;8;9;12;18;24 \} \)}{\( \{ 2;3;6 \} \)}{\( \{ 1;2;3;4;6;9;12 \} \)}{}{}}
\LANGcs{
\Question{
Urči množinu všech společných dělitelů čísel \( 24 \) a \( 18 \).}
{\( \{ 1;2;3;6 \} \)}{\( \{ 1;2;3;4;6;8;9;12;18;24 \} \)}{\( \{ 2;3;6 \} \)}{\( \{ 1;2;3;4;6;9;12 \} \)}{}{}}
\LANGsk{
\Question{
Urč množinu všetkých spoločných deliteľov čísel \( 24 \) a \( 18 \).}
{\( \{ 1;2;3;6 \} \)}{\( \{ 1;2;3;4;6;8;9;12;18;24 \} \)}{\( \{ 2;3;6 \} \)}{\( \{ 1;2;3;4;6;9;12 \} \)}{}{}}
\LANGes{
\Question{
Calcula el conjunto de todos los divisores comunes de los números  \( 24 \) y \( 18 \).}
{\( \{ 1;2;3;6 \} \)}{\( \{ 1;2;3;4;6;8;9;12;18;24 \} \)}{\( \{ 2;3;6 \} \)}{\( \{ 1;2;3;4;6;9;12 \} \)}{}{}}
\End

\Data\ID{1003055712}
\Updated{2020-07-15 03:07:12}
\Level{A}
\SubArea{Logic and Sets}
\LANGen{
\Question{
The set \( A \) is the set of all divisors of \( 12 \), the set \( B \) is the set of all divisors of \( 6 \). Which relation is correct?}
{\( B\subset A\)}{\( A\subset B\)}{\( A=B \)}{\( A\cup B=\{1;2;3;6;12\} \)}{}{}}
\LANGpl{
\Question{
Zbiór \( A \) to zbiór dzielników liczby  \( 12 \), zbiór \( B \) to zbiór dzielników liczby  \( 6 \). Zatem:}
{\( B\subset A\)}{\( A\subset B\)}{\( A=B \)}{\( A\cup B=\{1;2;3;6;12\} \)}{}{}}
\LANGcs{
\Question{
Množina \( A \) je množinou všech dělitelů \( 12 \), množina \( B \) je množinou všech dělitelů \( 6 \). Který vztah je správný?}
{\( B\subset A\)}{\( A\subset B\)}{\( A=B \)}{\( A\cup B=\{1;2;3;6;12\} \)}{}{}}
\LANGsk{
\Question{
Množina \( A \) je množinou všetkých deliteľov \( 12 \), množina \( B \) je množinou všetkých deliteľov \( 6 \). Ktorý vzťah je správny?}
{\( B\subset A\)}{\( A\subset B\)}{\( A=B \)}{\( A\cup B=\{1;2;3;6;12\} \)}{}{}}
\LANGes{
\Question{
El conjunto  \( A \) es el conjunto de todos los divisores de \( 12 \), el conjunto \( B \) es el conjunto de todos los divisores de \( 6 \). ¿Qué relación es correcta?}
{\( B\subset A\)}{\( A\subset B\)}{\( A=B \)}{\( A\cup B=\{1;2;3;6;12\} \)}{}{}}
\End

\Data\ID{1003055713}
\Updated{2020-07-15 03:07:12}
\Level{A}
\SubArea{Logic and Sets}
\LANGen{
\Question{
Choose the correct expression.}
{\( \{d\}\cup \{d\}\cup \{d\} = \{d\} \)}{\( \{d\}\cup \{d\}\cup \{d\} = \{ddd\} \)}{\( \{d\}\cup \{d\}\cup \{d\} = \{d;d\} \)}{\( \{d\}\cup \{d\}\cup \{d\} = \{d;d;d\} \)}{}{}}
\LANGpl{
\Question{
Zaznacz poprawne działanie:}
{\( \{d\}\cup \{d\}\cup \{d\} = \{d\} \)}{\( \{d\}\cup \{d\}\cup \{d\} = \{ddd\} \)}{\( \{d\}\cup \{d\}\cup \{d\} = \{d;d\} \)}{\( \{d\}\cup \{d\}\cup \{d\} = \{d;d;d\} \)}{}{}}
\LANGcs{
\Question{
Vyberte správný výraz.}
{\( \{d\}\cup \{d\}\cup \{d\} = \{d\} \)}{\( \{d\}\cup \{d\}\cup \{d\} = \{ddd\} \)}{\( \{d\}\cup \{d\}\cup \{d\} = \{d;d\} \)}{\( \{d\}\cup \{d\}\cup \{d\} = \{d;d;d\} \)}{}{}}
\LANGsk{
\Question{
Vyberte správny výraz.}
{\( \{d\}\cup \{d\}\cup \{d\} = \{d\} \)}{\( \{d\}\cup \{d\}\cup \{d\} = \{ddd\} \)}{\( \{d\}\cup \{d\}\cup \{d\} = \{d;d\} \)}{\( \{d\}\cup \{d\}\cup \{d\} = \{d;d;d\} \)}{}{}}
\LANGes{
\Question{
Elige la expresión correcta.}
{\( \{d\}\cup \{d\}\cup \{d\} = \{d\} \)}{\( \{d\}\cup \{d\}\cup \{d\} = \{ddd\} \)}{\( \{d\}\cup \{d\}\cup \{d\} = \{d;d\} \)}{\( \{d\}\cup \{d\}\cup \{d\} = \{d;d;d\} \)}{}{}}
\End

\Data\ID{1103055614}
\Updated{2022-06-30 02:06:40}
\Level{A}
\SubArea{Logic and Sets}
\IMG{1103055614.pdf}
\LANGen{
\Question{
Determine which of the following conditions for \( x \) defines the same set as is defined by the given diagram.}
{\( x\in(-3;1] \)}{\( x\in[-3;1) \)}{\( -3\leq x\leq1 \)}{\( x \leq 1 \) or \(x\geq -3\)}{}{}}
\LANGpl{
\Question{
Wskaż, który z poniższych warunków opisuje zbiór  zaznaczony na rysunku.}
{\( x\in(-3;1\rangle \)}{\( x\in\langle-3;1) \)}{\( -3\leq x\leq1 \)}{\( x \leq 1 \) lub \(x\geq -3\)}{}{}}
\LANGcs{
\Question{
Určete, která z následujících podmínek pro \( x \) definuje stejnou množinu, jaká je určena daným obrázkem.}
{\( x\in(-3;1\rangle \)}{\( x\in\langle-3;1) \)}{\( -3\leq x\leq1 \)}{\( x \leq 1 \) nebo \(x\geq -3\)}{}{}}
\LANGsk{
\Question{
Určte, ktorá z nasledujúcich podmienok pre \( x \) definuje rovnakú množinu, aká je určená daným obrázkom.}
{\( x\in(-3;1\rangle \)}{\( x\in\langle-3;1) \)}{\( -3\leq x\leq1 \)}{\( x \leq 1 \) alebo \(x\geq -3\)}{}{}}
\LANGes{
\Question{
Determina cuál de las siguientes condiciones para  \( x \) define el mismo conjunto  representado en el diagrama dado.}
{\( x\in(-3;1] \)}{\( x\in[-3;1) \)}{\( -3\leq x\leq1 \)}{\( x \leq 1 \) o \(x\geq -3\)}{}{}}
\End

\Data\ID{1103055703}
\Updated{2020-07-15 03:07:12}
\Level{A}
\SubArea{Logic and Sets}
\IMG{1103055703.pdf}
\LANGen{
\Question{
The diagram shows the sets \( A \) and \( B \). Find the union \( A\cup B \).}
{\( [ -3;4 ] \)}{\( ( -3;4 ) \)}{\( [ -1;2 ] \)}{\( ( -1;2 ) \)}{}{}}
\LANGpl{
\Question{
Sumą zbiorów \( A \) i \( B \) (\( A\cup B \)) jest:}
{\( \langle -3;4 \rangle \)}{\( ( -3;4 ) \)}{\( \langle -1;2 \rangle \)}{\( ( -1;2 ) \)}{}{}}
\LANGcs{
\Question{
Obrázek znázorňuje množiny\( A \) a \( B \). Najdi sjednocení \( A\cup B \).}
{\( \langle -3;4 \rangle \)}{\( ( -3;4 ) \)}{\( \langle -1;2 \rangle \)}{\( ( -1;2 ) \)}{}{}}
\LANGsk{
\Question{
Obrázok znázorňuje množiny\( A \) a \( B \). Nájdi zjednotenie \( A\cup B \).}
{\( \langle -3;4 \rangle \)}{\( ( -3;4 ) \)}{\( \langle -1;2 \rangle \)}{\( ( -1;2 ) \)}{}{}}
\LANGes{
\Question{
El diagrama muestra los conjuntos  \( A \) y \( B \). Calcula la unión \( A\cup B \).}
{\( [ -3;4 ] \)}{\( ( -3;4 ) \)}{\( [ -1;2 ] \)}{\( ( -1;2 ) \)}{}{}}
\End

\Data\ID{1103055714}
\Updated{2020-07-15 03:07:12}
\Level{A}
\SubArea{Logic and Sets}
\IMG{1103055714.pdf}
\LANGen{
\Question{
In the diagram below yellow colour represents}
{the difference of sets \( A \) and \( B \)}{the intersection of sets \( A \) and \( B \)}{the difference of sets \( B \) and \( A \)}{the union of sets \( A \) and \( B \)}{}{}}
\LANGpl{
\Question{
Na ilustracji poniżej kolorem żółtym zaznaczono:}
{różnicę zbiorów \( A \) i \( B \)}{iloczyn zbiorów \( A \) i \( B \)}{różnicę zbiorów \( B \) i \( A \)}{sumę zbiorów \( A \) i \( B \)}{}{}}
\LANGcs{
\Question{
Žlutá barva na obrázku znázorňuje}
{rozdíl množin \( A \) a \( B \)}{průnik množin \( A \) a \( B \)}{rozdíl množin \( B \) a \( A \)}{sjednocení množin \( A \) a \( B \)}{}{}}
\LANGsk{
\Question{
Žltá farba na obrázku znázorňuje}
{rozdiel množín \( A \) a \( B \)}{prienik množín \( A \) a \( B \)}{rozdiel množín \( B \) a \( A \)}{zjednotenie množín \( A \) a \( B \)}{}{}}
\LANGes{
\Question{
En el siguiente diagrama, el color amarillo representa}
{la diferencia de los conjuntos \( A \) y \( B \)}{la intersección de los conjuntos  \( A \) y \( B \)}{la diferencia de los conjuntos  \( B \) y \( A \)}{la unión de los conjuntos  \( A \) y \( B \)}{}{}}
\End

\Data\ID{2010000602}
\Updated{2021-12-01 04:12:54}
\Level{A}
\SubArea{Logic and Sets}
\LANGen{
\Question{
Given the sets \(A = [ -8;3 ]\) and \(B =(0;10)\) find the set difference \(A \setminus B\).}
{\( [ -8;0 ] \)}{\( [ -8;0 )\)}{\( ( 0;10)\)}{\( ( 0;3 ]\)}{}{}}
\LANGpl{
\Question{
Dane są zbiory \(A = \langle -8;3 \rangle\) i \(B =(0;10)\). Znajdź różnicę zbioru \(A \setminus B\).}
{\( \langle -8;0 \rangle \)}{\( \langle -8;0 )\)}{\( ( 0;10)\)}{\( ( 0;3 \rangle\)}{}{}}
\LANGcs{
\Question{
Určete rozdíl množin \(A \setminus B\), když \(A = \langle -8;3 \rangle\) a \(B =(0;10)\).}
{\( \langle -8;0 \rangle \)}{\( \langle -8;0 )\)}{\( ( 0;10)\)}{\( ( 0;3 \rangle\)}{}{}}
\LANGsk{
\Question{
Určte rozdiel množín \( A\setminus B \), ak
\(A = \langle -8;3 \rangle\), \(B =(0;10)\).}
{\( \langle -8;0 \rangle \)}{\( \langle -8;0 )\)}{\( ( 0;10)\)}{\( ( 0;3 \rangle\)}{}{}}
\LANGes{
\Question{
Dados los conjuntos   \(A = [ -8;3 ]\) y \(B =(0;10)\). Determina la diferencia de conjuntos \(A \setminus B\).}
{\( [ -8;0 ] \)}{\( [ -8;0 )\)}{\( ( 0;10)\)}{\( ( 0;3 ]\)}{}{}}
\End

\Data\ID{2010001403}
\Updated{2021-12-01 04:12:11}
\Level{A}
\SubArea{Logic and Sets}
\LANGen{
\Question{
Find the set difference \( A\setminus B \) for  \( A=\left\{x\in \mathbb{Z}\ \colon x^2=1\right\} \) and \( B=\{0;1;2;3\} \).}
{\( \{-1\} \)}{\( \{0;1;2;3\} \)}{\( \{0;2;3\} \)}{\( \emptyset\)}{}{}}
\LANGpl{
\Question{
Wyznacz różnicę zbiorów \( A\setminus B \), jeśli  \( A=\left\{x\in \mathbb{Z}\ \colon x^2=1\right\} \) i \( B=\{0;1;2;3\} \).}
{\( \{-1\} \)}{\( \{0;1;2;3\} \)}{\( \{0;2;3\} \)}{\( \emptyset\)}{}{}}
\LANGcs{
\Question{
Určete rozdíl množin \( A\setminus B \), jestliže  \( A=\left\{x\in \mathbb{Z}\ \colon x^2=1\right\} \) a \( B=\{0;1;2;3\} \).}
{\( \{-1\} \)}{\( \{0;1;2;3\} \)}{\( \{0;2;3\} \)}{\( \emptyset\)}{}{}}
\LANGsk{
\Question{
Určte rozdiel množín \( A\setminus B \) pre  \( A=\left\{x\in \mathbb{Z}\ \colon x^2=1\right\} \) a \( B=\{0;1;2;3\} \).}
{\( \{-1\} \)}{\( \{0;1;2;3\} \)}{\( \{0;2;3\} \)}{\( \emptyset\)}{}{}}
\LANGes{
\Question{
Determina la diferencia de conjuntos \( A\setminus B \) para  \( A=\left\{x\in \mathbb{Z}\ \colon x^2=1\right\} \) y \( B=\{0;1;2;3\} \).}
{\( \{-1\} \)}{\( \{0;1;2;3\} \)}{\( \{0;2;3\} \)}{\( \emptyset\)}{}{}}
\End

\Data\ID{2010001404}
\Updated{2021-12-01 04:12:16}
\Level{A}
\SubArea{Logic and Sets}
\LANGen{
\Question{
Find the intersection \( A\cap B' \) if \( A=(-\infty;4) \) and \(B=(-6;+\infty) \). (By \(B'\) the complement of the set \( B \) is denoted.)}
{\( (-\infty;-6 ] \)}{\( (-6;4) \)}{\( [ 4 ;+\infty) \)}{\( (-6;4 ] \)}{}{}}
\LANGpl{
\Question{
Wyznacz \( A\cap B' \) jeżeli \( A=(-\infty;4) \) i \(B=(-6;+\infty) \). (Przez \(B'\) oznaczamy dopełnienie zbioru \( B \).)}
{\( (-\infty;-6 \rangle \)}{\( (-6;4) \)}{\( \langle 4 ;+\infty) \)}{\( (-6;4 \rangle \)}{}{}}
\LANGcs{
\Question{
Najděte průnik množin \( A\cap B' \), jestliže \( A=(-\infty;4) \) a \(B=(-6;+\infty) \). (\(B'\) označuje doplněk množiny \( B \).)}
{\( (-\infty;-6 \rangle \)}{\( (-6;4) \)}{\( \langle 4 ;+\infty) \)}{\( (-6;4 \rangle \)}{}{}}
\LANGsk{
\Question{
Určte prienik množín \( A\cap B' \), ak \( A=(-\infty;4) \) a \(B=(-6;+\infty) \). (\(B'\) označujeme doplnok množiny \( B \).)}
{\( (-\infty;-6 \rangle \)}{\( (-6;4) \)}{\( \langle 4 ;+\infty) \)}{\( (-6;4 \rangle \)}{}{}}
\LANGes{
\Question{
Determina la intersección \( A\cap B' \) si \( A=(-\infty;4) \) y \(B=(-6;+\infty) \). \(B'\) es el complemento del conjunto \( B \).}
{\( (-\infty;-6 ] \)}{\( (-6;4) \)}{\( [ 4 ;+\infty) \)}{\( (-6;4 ] \)}{}{}}
\End

\Data\ID{9000086603}
\Updated{2020-07-15 03:07:37}
\Level{A}
\SubArea{Logic and Sets}
\LANGen{
\Question{
Determine the truth values of propositions \(a\) and \(b\) if you know that the compound proposition
\[ 
 \neg a \wedge b
\]
is true.}
{The statement \(a\) 
is false, \(b\) 
is true.}{Both statements are true.}{The statement \(a\) 
is true, \(b\) 
is false.}{Both statements are false.}{}{}}
\LANGpl{
\Question{
Oceń wartość logiczną zdań  \(a\) 
i  \(b\), jeśli wiadomo, że zdanie
\[ 
 \neg a \wedge b
\]
jest prawdziwe.}
{Zdanie \(a\) jest fałszywe, zdanie \(b\) jest prawdziwe.}{Oba zdania są prawdziwe.}{Zdanie \(a\) jest prawdziwe, zdanie \(b\) jest fałszywe.}{Oba zdania są fałszywe.}{}{}}
\LANGcs{
\Question{
Určete pravdivostní hodnoty výroků
\(a\) a
\(b\),
víte-li, že pravdivostní hodnota složeného výroku
\(\neg a \wedge b\) je
\(1\).}
{Pravdivostní hodnota \(a\) 
je \(0\), pravdivostní
hodnota \(b\) 
je \(1\).}{Pravdivostní hodnota \(a\) 
je \(1\), pravdivostní
hodnota \(b\) 
je \(1\).}{Pravdivostní hodnota \(a\) 
je \(1\), pravdivostní
hodnota \(b\) 
je \(0\).}{Pravdivostní hodnota \(a\) 
je \(0\), pravdivostní
hodnota \(b\) 
je \(0\).}{}{}}
\LANGsk{
\Question{
Určte pravdivostné hodnoty výrokov
\(a\) a
\(b\),
ak viete, že pravdivostná hodnota zloženého výroku
\(\neg a \wedge b\) je
\(1\).}
{Pravdivostná hodnota \(a\) 
je \(0\), pravdivostná
hodnota \(b\) 
je \(1\).}{Pravdivostná hodnota \(a\) 
je \(1\), pravdivostná
hodnota \(b\) 
je \(1\).}{Pravdivostná hodnota \(a\) 
je \(1\), pravdivostná
hodnota \(b\) 
je \(0\).}{Pravdivostná hodnota \(a\) 
je \(0\), pravdivostná
hodnota \(b\) 
je \(0\).}{}{}}
\LANGes{
\Question{
Determina los valores de verdad de las proposiciones  \(a\) y \(b\) b sabiendo que la proposición compuesta
\[ 
 \neg a \wedge b
\]
es verdadera.}
{La proposición \(a\) 
es falsa, \(b\) 
es verdadera.}{Ambas proposiciones son verdaderas.}{La proposición  \(a\) 
es verdadera, \(b\) 
es falsa.}{Ambas proposiciones son falsas.}{}{}}
\End

\Data\ID{9000148906}
\Updated{2020-07-15 03:07:37}
\Level{A}
\SubArea{Logic and Sets}
\LANGen{
\Question{
Each candidate in a tender is fluent in at least one out
of two required languages (English and French). There are
\(20\) candidates fluent
in English and \(14\) 
candidates fluent in French. From these amounts
\(10\) 
candidates are fluent in both languages. How many candidates are there in the
tender?}
{\(24\)}{\(34\)}{\(14\)}{\(44\)}{}{}}
\LANGpl{
\Question{
Każdy kandydat w ofercie jest biegły w przynajmniej jednym z dwóch wymaganych języków
 (język angielski, język francuski). 
\(20\) kandydatów jest biegłych w języku angielskim \(14\) 
w języku francuskim,
\(10\) 
posługuje się oboma językami. Ilu kandydatów jest w ofercie?}
{\(24\)}{\(34\)}{\(14\)}{\(44\)}{}{}}
\LANGcs{
\Question{
Ve skupině uchazečů o práci ovládá každý uchazeč alespoň jeden ze dvou jazyků (angličtinu nebo francouzštinu).
\(20\) uchazečů ovládá
angličtinu a \(14\) 
uchazečů ovládá francouzštinu. Přitom
\(10\) 
uchazečů ovládá oba jazyky. Kolik uchazečů je na konkurzu?}
{\(24\)}{\(34\)}{\(14\)}{\(44\)}{}{}}
\LANGsk{
\Question{
V skupine uchádzačov o prácu ovláda každý uchádzač plynule aspoň jeden z dvoch cudzích jazykov (angličtinu alebo francúzštinu).
\(20\) uchádzačov ovláda plynule anglický jazyk a 
 \(14\) 
uchádzačov ovláda plynule francúzsky jazyk. Pritom
\(10\) uchádzačov ovláda obidva jazyky.
Koľko uchádzačov je na konkurze?}
{\(24\)}{\(34\)}{\(14\)}{\(44\)}{}{}}
\LANGes{
\Question{
Cada candidato de un concurso habla con fluidez al menos uno de los dos idiomas requeridos (inglés y francés). Hay
\(20\) candidatos con fluidez en inglés y  \(14\) 
candidatos con fluidez en francés. De estas cantidades
\(10\) 
candidatos dominan ambos idiomas. ¿Cuántos candidatos hay en el concurso?}
{\(24\)}{\(34\)}{\(14\)}{\(44\)}{}{}}
\End

\Data\ID{1003034301}
\Updated{2022-08-25 10:08:00}
\Level{B}
\SubArea{Logic and Sets}
\LANGen{
\Question{
Given the sets \( A=(-1;+\infty) \) and \( B=[-3;6) \), find the intersection  \( A'\cap B \).}
{\( [-3;-1] \)}{\( [-1;6) \)}{\( [-3;-1) \)}{\( (-3;-1) \)}{}{}}
\LANGpl{
\Question{
Jeżeli  \( A=(-1;+\infty) \) i \( B=\langle-3;6) \), to  \( A'\cap B \) jest równe:}
{\( \langle-3;-1\rangle \)}{\( \langle-1;6) \)}{\( \langle-3;-1) \)}{\( (-3;-1) \)}{}{}}
\LANGcs{
\Question{
Mějme množiny \( A=(-1;+\infty) \) a \( B=\langle-3;6) \), Určete průnik  \( A'\cap B \).}
{\( \langle-3;-1\rangle \)}{\( \langle-1;6) \)}{\( \langle-3;-1) \)}{\( (-3;-1) \)}{}{}}
\LANGsk{
\Question{
Máme dané množiny \( A=(-1;+\infty) \) a \( B=\langle-3;6) \), najdi prienik  \( A'\cap B \).}
{\( \langle-3;-1\rangle \)}{\( \langle-1;6) \)}{\( \langle-3;-1) \)}{\( (-3;-1) \)}{}{}}
\LANGes{
\Question{
Dados los conjuntos \( A=(-1;+\infty) \) y \( B=[-3;6) \), calcula la intersección \( A'\cap B \).}
{\( [-3;-1] \)}{\( [-1;6) \)}{\( [-3;-1) \)}{\( (-3;-1) \)}{}{}}
\End

\Data\ID{1003034302}
\Updated{2022-08-25 10:08:00}
\Level{B}
\SubArea{Logic and Sets}
\LANGen{
\Question{
Given the sets \( A=(-6;2] \), \( B=[-2;8) \) and \( C=(0;3] \), determine \( (A\cap B)\setminus C \).}
{\( [-2;0] \)}{\( [-2;0) \)}{\( [-6;8) \)}{\( (-6;0)\cup(3;8) \)}{}{}}
\LANGpl{
\Question{
Dane są zbiory \( A=(-6;2\rangle \), \( B=\langle-2;8) \) i \( C=(0;3\rangle \), wynikiem działań \( (A\cap B)\setminus C \) jest zbiór:}
{\( \langle-2;0\rangle \)}{\( \langle-2;0) \)}{\( \langle-6;8) \)}{\( (-6;0)\cup(3;8) \)}{}{}}
\LANGcs{
\Question{
Mějme množiny \( A=(-6;2\rangle \), \( B=\langle-2;8) \) a \( C=(0;3\rangle \), určete \( (A\cap B)\setminus C \).}
{\( \langle-2;0\rangle \)}{\( \langle-2;0) \)}{\( \langle-6;8) \)}{\( (-6;0)\cup(3;8) \)}{}{}}
\LANGsk{
\Question{
Máme dané množiny \( A=(-6;2\rangle \), \( B=\langle-2;8) \) a \( C=(0;3\rangle \), urč \( (A\cap B)\setminus C \).}
{\( \langle-2;0\rangle \)}{\( \langle-2;0) \)}{\( \langle-6;8) \)}{\( (-6;0)\cup(3;8) \)}{}{}}
\LANGes{
\Question{
Dados los conjuntos \( A=(-6;2] \), \( B=[-2;8) \) y \( C=(0;3] \), determina \( (A\cap B)\setminus C \).}
{\( [-2;0] \)}{\( [-2;0) \)}{\( [-6;8) \)}{\( (-6;0)\cup(3;8) \)}{}{}}
\End

\Data\ID{1003034303}
\Updated{2022-08-25 10:08:00}
\Level{B}
\SubArea{Logic and Sets}
\LANGen{
\Question{
Let \( A=(-\infty;2] \), \( B=[-4;3) \) and \( C=(-\infty;2] \). Determine \( (A\cup B)\setminus C \).}
{\( (2;3) \)}{\( [2;3) \)}{\( (-\infty;-4) \)}{\( (-\infty;3) \)}{}{}}
\LANGpl{
\Question{
Dane są zbiory \( A=(-\infty;2\rangle \), \( B=\langle-4;3) \) i \( C=(-\infty;2\rangle \). Wynikiem działań \( (A\cup B)\setminus C \) jest zbiór:}
{\( (2;3) \)}{\( \langle2;3) \)}{\( (-\infty;-4) \)}{\( (-\infty;3) \)}{}{}}
\LANGcs{
\Question{
Nechť \( A=(-\infty;2\rangle \), \( B=\langle-4;3) \) a \( C=(-\infty;2\rangle \). Určete \( (A\cup B)\setminus C \).}
{\( (2;3) \)}{\( \langle2;3) \)}{\( (-\infty;-4) \)}{\( (-\infty;3) \)}{}{}}
\LANGsk{
\Question{
Nech \( A=(-\infty;2\rangle \), \( B=\langle-4;3) \) a \( C=(-\infty;2\rangle \). Urč \( (A\cup B)\setminus C \).}
{\( (2;3) \)}{\( \langle2;3) \)}{\( (-\infty;-4) \)}{\( (-\infty;3) \)}{}{}}
\LANGes{
\Question{
Sean \( A=(-\infty;2] \), \( B=[-4;3) \) y \( C=(-\infty;2] \). Determina \( (A\cup B)\setminus C \).}
{\( (2;3) \)}{\( [2;3) \)}{\( (-\infty;-4) \)}{\( (-\infty;3) \)}{}{}}
\End

\Data\ID{1003034304}
\Updated{2022-08-25 10:08:00}
\Level{B}
\SubArea{Logic and Sets}
\LANGen{
\Question{
Given the sets \( A=(-\infty;-5] \), \( B=[-5;0) \) and \( C=(-7;-1) \), find \( A\setminus(B\cap C) \).}
{\( (-\infty;-5) \)}{\( (-\infty;-5] \)}{\( [-5;-1) \)}{\( (-\infty;0) \)}{}{}}
\LANGpl{
\Question{
Dane są zbiory \( A=(-\infty;-5\rangle \), \( B=\langle-5;0) \) i \( C=(-7;-1) \), Wynikiem działań  \( A\setminus(B\cap C) \) jest zbiór:}
{\( (-\infty;-5) \)}{\( (-\infty;-5\rangle \)}{\( \langle-5;-1) \)}{\( (-\infty;0) \)}{}{}}
\LANGcs{
\Question{
Máme dané množiny \( A=(-\infty;-5\rangle \), \( B=\langle-5;0) \) a \( C=(-7;-1) \), najděte \( A\setminus(B\cap C) \).}
{\( (-\infty;-5) \)}{\( (-\infty;-5\rangle \)}{\( \langle-5;-1) \)}{\( (-\infty;0) \)}{}{}}
\LANGsk{
\Question{
Máme dané množiny \( A=(-\infty;-5\rangle \), \( B=\langle-5;0) \) a \( C=(-7;-1) \), nájdi \( A\setminus(B\cap C) \).}
{\( (-\infty;-5) \)}{\( (-\infty;-5\rangle \)}{\( \langle-5;-1) \)}{\( (-\infty;0) \)}{}{}}
\LANGes{
\Question{
Dados los conjuntos \( A=(-\infty;-5] \), \( B=[-5;0) \) y \( C=(-7;-1) \), calcula \( A\setminus(B\cap C) \).}
{\( (-\infty;-5) \)}{\( (-\infty;-5] \)}{\( [-5;-1) \)}{\( (-\infty;0) \)}{}{}}
\End

\Data\ID{1003034305}
\Updated{2022-08-25 10:08:00}
\Level{B}
\SubArea{Logic and Sets}
\LANGen{
\Question{
Let \( A=(-\infty;-4)\cup(2;+\infty) \) and \( B = [ 2;8 ) \). Give the intersection \( A'\cap B \).}
{\( \{2\} \)}{\( \emptyset \)}{\( (2;8) \)}{\( [2;8) \)}{}{}}
\LANGpl{
\Question{
Jeżeli \( A=(-\infty;-4)\cup(2;+\infty) \) i \( B = \langle 2;8 ) \). To \( A'\cap B \) jest równe:}
{\( \{2\} \)}{\( \emptyset \)}{\( (2;8) \)}{\( \langle2;8) \)}{}{}}
\LANGcs{
\Question{
Nechť \( A=(-\infty;-4)\cup(2;+\infty) \) a \( B = \langle 2;8 ) \). Určete průnik \( A'\cap B \).}
{\( \{2\} \)}{\( \emptyset \)}{\( (2;8) \)}{\( \langle2;8) \)}{}{}}
\LANGsk{
\Question{
Nech \( A=(-\infty;-4)\cup(2;+\infty) \) a \( B = \langle 2;8 ) \). Urč prienik \( A'\cap B \).}
{\( \{2\} \)}{\( \emptyset \)}{\( (2;8) \)}{\( \langle2;8) \)}{}{}}
\LANGes{
\Question{
Sean \( A=(-\infty;-4)\cup(2;+\infty) \) y \( B = [ 2;8 ) \). Calcula la intersección \( A'\cap B \).}
{\( \{2\} \)}{\( \emptyset \)}{\( (2;8) \)}{\( [2;8) \)}{}{}}
\End

\Data\ID{1003034306}
\Updated{2022-08-25 10:08:03}
\Level{B}
\SubArea{Logic and Sets}
\LANGen{
\Question{
Choose another correct description of the set \( A=(-3;6)\cap(-2;+\infty)\cap\mathbb{Z} \).}
{\( A=\{-1;0;1;2;3;4;5\} \)}{\( A=\{-2,-1;0;1;2;3;4;5\} \)}{\( A=\{0;1;2;3;4;5\} \)}{\( A=(-2;6) \)}{}{}}
\LANGpl{
\Question{
Zbiór \( A=(-3;6)\cap(-2;+\infty)\cap\mathbb{Z} \)  jest równy:}
{\( A=\{-1;0;1;2;3;4;5\} \)}{\( A=\{-2,-1;0;1;2;3;4;5\} \)}{\( A=\{0;1;2;3;4;5\} \)}{\( A=(-2;6) \)}{}{}}
\LANGcs{
\Question{
Vyberte alternativní zápis množiny \( A=(-3;6)\cap(-2;+\infty)\cap\mathbb{Z} \).}
{\( A=\{-1;0;1;2;3;4;5\} \)}{\( A=\{-2,-1;0;1;2;3;4;5\} \)}{\( A=\{0;1;2;3;4;5\} \)}{\( A=(-2;6) \)}{}{}}
\LANGsk{
\Question{
Vyber ďalší správny zápis množiny \( A=(-3;6)\cap(-2;+\infty)\cap\mathbb{Z} \).}
{\( A=\{-1;0;1;2;3;4;5\} \)}{\( A=\{-2,-1;0;1;2;3;4;5\} \)}{\( A=\{0;1;2;3;4;5\} \)}{\( A=(-2;6) \)}{}{}}
\LANGes{
\Question{
Elige otra descripción correcta del conjunto \( A=(-3;6)\cap(-2;+\infty)\cap\mathbb{Z} \).}
{\( A=\{-1;0;1;2;3;4;5\} \)}{\( A=\{-2,-1;0;1;2;3;4;5\} \)}{\( A=\{0;1;2;3;4;5\} \)}{\( A=(-2;6) \)}{}{}}
\End

\Data\ID{1003034307}
\Updated{2022-08-25 10:08:03}
\Level{B}
\SubArea{Logic and Sets}
\LANGen{
\Question{
Choose the correct description of the set \( X=(22;31)\cap\mathbb{N} \).}
{It is a set with eight elements.}{It is a set with ten elements.}{It is an open interval.}{It is a set of all natural numbers.}{}{}}
\LANGpl{
\Question{
Zbiór  \( X=(22;31)\cap\mathbb{N} \) jest:}
{Zbiorem ośmioelementowym.}{Zbiorem dziesięcioelementowym.}{Przedziałem obustronnie otwartym.}{Zbiorem liczb naturalnych.}{}{}}
\LANGcs{
\Question{
Vyberte správný popis množiny \( X=(22;31)\cap\mathbb{N} \).}
{Množina má osm prvků.}{Množina má deset prvků.}{Je to otevřený interval.}{Je to množina všech přirozených čísel.}{}{}}
\LANGsk{
\Question{
Vyber správny zápis množiny \( X=(22;31)\cap\mathbb{N} \).}
{Množina má osem prvkov.}{Množina má desať prvkov.}{Je to otvorený interval.}{Je to množina všetkých prirodzených čísel.}{}{}}
\LANGes{
\Question{
Elige otra afirmación correcta del conjunto \( X=(22;31)\cap\mathbb{N} \).}
{Es un conjunto de ocho elementos.}{Es un conjunto de diez elementos.}{Es un conjunto abierto.}{Es el conjunto de los números naturales.}{}{}}
\End

\Data\ID{1003034308}
\Updated{2022-08-25 10:08:03}
\Level{B}
\SubArea{Logic and Sets}
\LANGen{
\Question{
Choose the correct description of the set \( X=[ -10;100 ]\cap\mathbb{Z} \).}
{It is a finite set.}{It is a closed interval.}{It is a set with \( 110 \) elements.}{It is a set with even number of elements.}{}{}}
\LANGpl{
\Question{
Zbiór \( X=\langle -10;100 \rangle\cap\mathbb{Z} \) jest:}
{Zbiorem o skończonej liczbie elementów.}{Przedziałem obustronnie domkniętym.}{Zbiorem, który ma \( 110 \) elementów.}{Zbiorem o parzystej liczbie elementów.}{}{}}
\LANGcs{
\Question{
Označte pravdivé tvrzení o množině \( X=\langle -10;100 \rangle\cap\mathbb{Z} \).}
{\( X \) je konečná množina.}{\( X \) je uzavřený interval.}{\( X \) obsahuje \( 110 \) prvků.}{Všechny prvky množiny \( X \) jsou sudá čísla.}{}{}}
\LANGsk{
\Question{
Označte pravdivé tvrdenie o množine \( X=\langle -10;100 \rangle\cap\mathbb{Z} \).}
{\( X \) je konečná množina.}{\( X \) je uzavretý interval.}{\( X \)  obsahuje \( 110 \) prvkov.}{Všetky prvky množiny \( X \) sú párne čísla.}{}{}}
\LANGes{
\Question{
Elige la descripción correcta del conjunto \( X=[ -10;100 ]\cap\mathbb{Z} \).}
{Es un conjunto finito.}{Es un intervalo cerrado.}{Es un conjunto de \( 110 \) elementos.}{Es un conjunto con un número par de elementos.}{}{}}
\End

\Data\ID{1003034309}
\Updated{2022-08-25 10:08:03}
\Level{B}
\SubArea{Logic and Sets}
\LANGen{
\Question{
Choose another correct description of the set   \( X=(2;14)\cap\{2;4;6;8;10;12;14\} \).}
{\( X=\{4;6;8;10;12\} \)}{\( X=\{2;4;6;8;10;12;14\} \)}{\( X=(2;14) \)}{\( X=[2;14] \)}{}{}}
\LANGpl{
\Question{
Zbiór  \( X=(2;14)\cap\{2;4;6;8;10;12;14\} \) jest równy:}
{\( X=\{4;6;8;10;12\} \)}{\( X=\{2;4;6;8;10;12;14\} \)}{\( X=(2;14) \)}{\( X=\langle2;14\rangle \)}{}{}}
\LANGcs{
\Question{
Vyberte alternativní zápis množiny   \( X=(2;14)\cap\{2;4;6;8;10;12;14\} \).}
{\( X=\{4;6;8;10;12\} \)}{\( X=\{2;4;6;8;10;12;14\} \)}{\( X=(2;14) \)}{\( X=\langle2;14\rangle \)}{}{}}
\LANGsk{
\Question{
Vyber ďalší správny zápis množiny   \( X=(2;14)\cap\{2;4;6;8;10;12;14\} \).}
{\( X=\{4;6;8;10;12\} \)}{\( X=\{2;4;6;8;10;12;14\} \)}{\( X=(2;14) \)}{\( X=\langle2;14\rangle \)}{}{}}
\LANGes{
\Question{
Elige otra descripción correcta del conjunto \( X=(2;14)\cap\{2;4;6;8;10;12;14\} \).}
{\( X=\{4;6;8;10;12\} \)}{\( X=\{2;4;6;8;10;12;14\} \)}{\( X=(2;14) \)}{\( X=[2;14] \)}{}{}}
\End

\Data\ID{1003034310}
\Updated{2022-08-25 10:08:03}
\Level{B}
\SubArea{Logic and Sets}
\LANGen{
\Question{
Find the union  \( A'\cup B \) for \( A=(-1;7) \) and \( B=[-3;+\infty) \).}
{\( \mathbb{R} \)}{\( [ -3;+\infty) \)}{\( [-3;-1]\cup[7;+\infty) \)}{\( (-3;-1)\cup(7;+\infty) \)}{}{}}
\LANGpl{
\Question{
Zbiór  \( A'\cup B \) dla \( A=(-1;7) \) i \( B=\langle-3;+\infty) \) jest równy:}
{\( \mathbb{R} \)}{\( \langle -3;+\infty) \)}{\( \langle-3;-1\rangle\cup\langle7;+\infty) \)}{\( (-3;-1)\cup(7;+\infty) \)}{}{}}
\LANGcs{
\Question{
Najděte sjednocení  \( A'\cup B \), jestliže \( A=(-1;7) \) a \( B=\langle-3;+\infty) \).}
{\( \mathbb{R} \)}{\( \langle -3;+\infty) \)}{\( \langle-3;-1\rangle\cup\langle7;+\infty) \)}{\( (-3;-1)\cup(7;+\infty) \)}{}{}}
\LANGsk{
\Question{
Nájdi zjednotenie  \( A'\cup B \), ak \( A=(-1;7) \) a \( B=\langle-3;+\infty) \).}
{\( \mathbb{R} \)}{\( \langle -3;+\infty) \)}{\( \langle-3;-1\rangle\cup\langle7;+\infty) \)}{\( (-3;-1)\cup(7;+\infty) \)}{}{}}
\LANGes{
\Question{
Calcula la unión \( A'\cup B \) para \( A=(-1;7) \) y \( B=[-3;+\infty) \).}
{\( \mathbb{R} \)}{\( [ -3;+\infty) \)}{\( [-3;-1]\cup[7;+\infty) \)}{\( (-3;-1)\cup(7;+\infty) \)}{}{}}
\End

\Data\ID{1003034311}
\Updated{2022-08-25 10:08:03}
\Level{B}
\SubArea{Logic and Sets}
\LANGen{
\Question{
Find the set difference  \( A'\setminus B' \) for \( A=(-1;2] \) and \( B=[ -3;4) \).}
{\( [-3;-1]\cup(2;4) \)}{\( [-3;-1]\cup[2;4] \)}{\( (-3;-1)\cup[2;4] \)}{\( (-3;-1]\cup(2;4] \)}{}{}}
\LANGpl{
\Question{
Zbiór  \( A'\setminus B' \) dla \( A=(-1;2\rangle \) i \( B=\langle -3;4) \) jest równy:}
{\( \langle-3;-1\rangle\cup(2;4) \)}{\( \langle-3;-1\rangle\cup\langle2;4\rangle \)}{\( (-3;-1)\cup\langle2;4\rangle \)}{\( (-3;-1\rangle\cup(2;4\rangle \)}{}{}}
\LANGcs{
\Question{
Najděte rozdíl množin  \( A'\setminus B' \), jestliže \( A=(-1;2\rangle \) a \( B=\langle -3;4) \).}
{\( \langle-3;-1\rangle\cup(2;4) \)}{\( \langle-3;-1\rangle\cup\langle2;4\rangle \)}{\( (-3;-1)\cup\langle2;4\rangle \)}{\( (-3;-1\rangle\cup(2;4\rangle \)}{}{}}
\LANGsk{
\Question{
Nájdi rozdiel množín  \( A'\setminus B' \), ak \( A=(-1;2\rangle \) a \( B=\langle -3;4) \).}
{\( \langle-3;-1\rangle\cup(2;4) \)}{\( \langle-3;-1\rangle\cup\langle2;4\rangle \)}{\( (-3;-1)\cup\langle2;4\rangle \)}{\( (-3;-1\rangle\cup(2;4\rangle \)}{}{}}
\LANGes{
\Question{
Calcula la diferencia de conjuntos  \( A'\setminus B' \) para \( A=(-1;2] \) y \( B=[ -3;4) \).}
{\( [-3;-1]\cup(2;4) \)}{\( [-3;-1]\cup[2;4] \)}{\( (-3;-1)\cup[2;4] \)}{\( (-3;-1]\cup(2;4] \)}{}{}}
\End

\Data\ID{1003034312}
\Updated{2022-08-25 10:08:03}
\Level{B}
\SubArea{Logic and Sets}
\LANGen{
\Question{
Choose another correct description of the set  \( A=\left((-3;8)\setminus[5;9)\right)\cap\mathbb{N} \).}
{\( A=\{1;2;3;4\} \)}{\( A=\{-3;-2;-1;0;1;2;3;4\} \)}{\( A=\{0;1;2;3;4;5\} \)}{\( A=\{1;2;3;4;5\} \)}{}{}}
\LANGpl{
\Question{
Elementami zbioru  \( A=\left((-3;8)\setminus\langle5;9)\right)\cap\mathbb{N} \) są:}
{\( A=\{1;2;3;4\} \)}{\( A=\{-3;-2;-1;0;1;2;3;4\} \)}{\( A=\{0;1;2;3;4;5\} \)}{\( A=\{1;2;3;4;5\} \)}{}{}}
\LANGcs{
\Question{
Vyberte alternativní zápis množiny \( A=\left((-3;8)\setminus\langle5;9)\right)\cap\mathbb{N} \).}
{\( A=\{1;2;3;4\} \)}{\( A=\{-3;-2;-1;0;1;2;3;4\} \)}{\( A=\{0;1;2;3;4;5\} \)}{\( A=\{1;2;3;4;5\} \)}{}{}}
\LANGsk{
\Question{
Vyber ďalší správny zápis množiny \( A=\left((-3;8)\setminus\langle5;9)\right)\cap\mathbb{N} \).}
{\( A=\{1;2;3;4\} \)}{\( A=\{-3;-2;-1;0;1;2;3;4\} \)}{\( A=\{0;1;2;3;4;5\} \)}{\( A=\{1;2;3;4;5\} \)}{}{}}
\LANGes{
\Question{
Elige otra descripción correcta del conjunto  \( A=\left((-3;8)\setminus[5;9)\right)\cap\mathbb{N} \).}
{\( A=\{1;2;3;4\} \)}{\( A=\{-3;-2;-1;0;1;2;3;4\} \)}{\( A=\{0;1;2;3;4;5\} \)}{\( A=\{1;2;3;4;5\} \)}{}{}}
\End

\Data\ID{1003034313}
\Updated{2022-08-25 10:08:03}
\Level{B}
\SubArea{Logic and Sets}
\LANGen{
\Question{
Find the set \( (B\setminus A) \cap C \) for \( A=(-10;2] \), \( B=[-2;8) \) and \( C=(0;4] \).}
{\( (2;4] \)}{\( (-10;0] \)}{\( [0;4] \)}{\( (8;4] \)}{}{}}
\LANGpl{
\Question{
Zbiór \( (B\setminus A) \cap C \) dla \( A=(-10;2\rangle \), \( B=\langle-2;8) \) i \( C=(0;4\rangle \) jest równy:}
{\( (2;4\rangle \)}{\( (-10;0\rangle \)}{\( \langle0;4\rangle \)}{\( (8;4\rangle \)}{}{}}
\LANGcs{
\Question{
Najděte množinu \( (B\setminus A) \cap C \), jestliže \( A=(-10;2\rangle \), \( B=\langle-2;8) \) a \( C=(0;4\rangle \).}
{\( (2;4\rangle \)}{\( (-10;0\rangle \)}{\( \langle0;4\rangle \)}{\( (8;4\rangle \)}{}{}}
\LANGsk{
\Question{
Nájdi množinu \( (B\setminus A) \cap C \), ak \( A=(-10;2\rangle \), \( B=\langle-2;8) \) a \( C=(0;4\rangle \).}
{\( (2;4\rangle \)}{\( (-10;0\rangle \)}{\( \langle0;4\rangle \)}{\( (8;4\rangle \)}{}{}}
\LANGes{
\Question{
Calcula el conjunto \( (B\setminus A) \cap C \) para \( A=(-10;2] \), \( B=[-2;8) \) y \( C=(0;4] \).}
{\( (2;4] \)}{\( (-10;0] \)}{\( [0;4] \)}{\( (8;4] \)}{}{}}
\End

\Data\ID{1003055605}
\Updated{2020-07-15 03:07:12}
\Level{B}
\SubArea{Logic and Sets}
\LANGen{
\Question{
Which of the following expressions is a proposition?}
{\( 5^2=25 \)}{\( 10^2 \)}{\( (a-b)^2 - (a+b)^2 \)}{Is it raining today?}{}{}}
\LANGpl{
\Question{
Które z poniższych wyrażeń jest zdaniem logicznym?}
{\( 5^2=25 \)}{\( 10^2 \)}{\( (a-b)^2 - (a+b)^2 \)}{Czy dzisiaj pada deszcz?}{}{}}
\LANGcs{
\Question{
Které z následujících vyjádření je výrok?}
{\( 5^2=25 \)}{\( 10^2 \)}{\( (a-b)^2 - (a+b)^2 \)}{Prší dnes?}{}{}}
\LANGsk{
\Question{
Ktoré z následujúcich vyjadrení je výrok?}
{\( 5^2=25 \)}{\( 10^2 \)}{\( (a-b)^2 - (a+b)^2 \)}{Prší dnes?}{}{}}
\LANGes{
\Question{
¿Cuál de las siguientes expresiones es una proposición lógica?}
{\( 5^2=25 \)}{\( 10^2 \)}{\( (a-b)^2 - (a+b)^2 \)}{¿Está lloviendo hoy?}{}{}}
\End

\Data\ID{1003055606}
\Updated{2020-07-15 03:07:12}
\Level{B}
\SubArea{Logic and Sets}
\LANGen{
\Question{
Identify a true proposition.}
{\( (2 < 5)\wedge(4 < 5) \)}{\( (0 < -1)\vee (8 >10) \)}{\( (3 < 5) \Rightarrow (8\leq -10) \)}{\( (5 < 10) \Leftrightarrow ( 8\div3=5) \)}{}{}}
\LANGpl{
\Question{
Które zdanie jest prawdziwe?}
{\( (2 < 5)\wedge(4 < 5) \)}{\( (0 < -1)\vee (8 >10) \)}{\( (3 < 5) \Rightarrow (8\leq -10) \)}{\( (5 < 10) \Leftrightarrow ( 8\div3=5) \)}{}{}}
\LANGcs{
\Question{
Urči pravdivý výrok.}
{\( (2 < 5)\wedge(4 < 5) \)}{\( (0 < -1)\vee (8 >10) \)}{\( (3 < 5) \Rightarrow (8\leq -10) \)}{\( (5 < 10) \Leftrightarrow ( 8\div3=5) \)}{}{}}
\LANGsk{
\Question{
Urč pravdivý výrok.}
{\( (2 < 5)\wedge(4 < 5) \)}{\( (0 < -1)\vee (8 >10) \)}{\( (3 < 5) \Rightarrow (8\leq -10) \)}{\( (5 < 10) \Leftrightarrow ( 8\div3=5) \)}{}{}}
\LANGes{
\Question{
Identifica una proposición lógica verdadera.}
{\( (2 < 5)\wedge(4 < 5) \)}{\( (0 < -1)\vee (8 >10) \)}{\( (3 < 5) \Rightarrow (8\leq -10) \)}{\( (5 < 10) \Leftrightarrow ( 8\div3=5) \)}{}{}}
\End

\Data\ID{1003055607}
\Updated{2020-07-15 03:07:12}
\Level{B}
\SubArea{Logic and Sets}
\LANGen{
\Question{
Find a false statement.}
{\( (5 < -10) \vee (4 < 3) \)}{\( (3\in\mathbb{N})\Leftrightarrow (3\in\mathbb{Z} ) \)}{\( (2 > 0)\vee(3=5) \)}{\( \left(8^2 = 16 \right) \Rightarrow \left(8^2 = 15 \right) \)}{}{}}
\LANGpl{
\Question{
Wskaż zdanie fałszywe.}
{\( (5 < -10) \vee (4 < 3) \)}{\( (3\in\mathbb{N})\Leftrightarrow (3\in\mathbb{Z} ) \)}{\( (2 > 0)\vee(3=5) \)}{\( \left(8^2 = 16 \right) \Rightarrow \left(8^2 = 15 \right) \)}{}{}}
\LANGcs{
\Question{
Urči nepravdivý výrok.}
{\( (5 < -10) \vee (4 < 3) \)}{\( (3\in\mathbb{N})\Leftrightarrow (3\in\mathbb{Z} ) \)}{\( (2 > 0)\vee(3=5) \)}{\( \left(8^2 = 16 \right) \Rightarrow \left(8^2 = 15 \right) \)}{}{}}
\LANGsk{
\Question{
Urč nepravdivý výrok.}
{\( (5 < -10) \vee (4 < 3) \)}{\( (3\in\mathbb{N})\Leftrightarrow (3\in\mathbb{Z} ) \)}{\( (2 > 0)\vee(3=5) \)}{\( \left(8^2 = 16 \right) \Rightarrow \left(8^2 = 15 \right) \)}{}{}}
\LANGes{
\Question{
Encuentra una proposición falsa.}
{\( (5 < -10) \vee (4 < 3) \)}{\( (3\in\mathbb{N})\Leftrightarrow (3\in\mathbb{Z} ) \)}{\( (2 > 0)\vee(3=5) \)}{\( \left(8^2 = 16 \right) \Rightarrow \left(8^2 = 15 \right) \)}{}{}}
\End

\Data\ID{1003055608}
\Updated{2022-06-30 02:06:17}
\Level{B}
\SubArea{Logic and Sets}
\LANGen{
\Question{
Determine the truth value of the statement \( \forall x\in\mathbb{R}\colon x^2+1>0 \).}
{It is a true statement.}{It is a false statement.}{It is not a statement.}{It is impossible to determine whether it is a true or a false statement.}{}{}}
\LANGpl{
\Question{
Oceń wartość logiczną zdania \( \forall x\in\mathbb{R}\colon x^2+1>0 \).}
{Jest zdaniem logicznym prawdziwym.}{Jest zdaniem logicznym fałszywym.}{Nie jest zdaniem logicznym.}{Nie można ocenić wartości logicznej.}{}{}}
\LANGcs{
\Question{
Rozhodněte, zda je výraz 
\[ \forall x\in\mathbb{R}\colon x^2+1>0 \]
výrokem, a v případě, že ano, určete jeho pravdivostní hodnotu.}
{Je to pravdivý výrok.}{Je to nepravdivý výrok.}{Není to výrok.}{Je nemožné určit, zda jde o pravdivý, nebo nepravdivý výrok.}{}{}}
\LANGsk{
\Question{
Urč pravdivú hodnotu výrazu \( \forall x\in\mathbb{R}\colon x^2+1>0 \).}
{Je to pravdivý výrok.}{Je to nepravdivý výrok.}{Nie je to výrok.}{Je nemožné určiť, či ide o pravdivý alebo nepravdivý výrok.}{}{}}
\LANGes{
\Question{
Determina el valor de verdad de la proposición lógica \( \forall x\in\mathbb{R}\colon x^2+1>0 \).}
{Es una proposición lógica verdadera.}{Es una proposición lógica falsa.}{No es una proposición lógica.}{Es imposible determinar si es una proposición lógica verdadera o falsa.}{}{}}
\End

\Data\ID{1003055609}
\Updated{2022-11-01 09:11:12}
\Level{B}
\SubArea{Logic and Sets}
\LANGen{
\Question{
Determine the truth value of the statement \( \exists k\in\mathbb{Z}\colon k^2 < 0\).}
{It is a false statement.}{It is a true statement.}{It is not a statement.}{It is impossible to determine whether it is a true or a false statement.}{}{}}
\LANGpl{
\Question{
Oceń wartość logiczną zdania \( \exists k\in\mathbb{Z}\colon k^2 < 0\).}
{Jest zdaniem logicznym fałszywym.}{Jest zdaniem logicznym prawdziwym.}{Nie jest zdaniem logicznym.}{Nie można ocenić wartości logicznej.}{}{}}
\LANGcs{
\Question{
Rozhodněte, zda je výraz 
\[ \exists k\in\mathbb{Z}\colon k^2 < 0\]
výrokem, a v případě, že ano, určete jeho pravdivostní hodnotu.}
{Je to nepravdivý výrok.}{Je to pravdivý výrok.}{Není to výrok.}{Je nemožné určit, zda jde o pravdivý, nebo nepravdivý výrok.}{}{}}
\LANGsk{
\Question{
Urč pravdivú hodnotu výroku \( \exists k\in\mathbb{Z}\colon k^2 < 0\).}
{Je to nepravdivý výrok.}{Je to pravdivý výrok.}{Nie je to výrok.}{Je nemožné určiť, či ide o pravdivý alebo nepravdivý výrok.}{}{}}
\LANGes{
\Question{
Determina el valor de verdad de la proposición lógica \( \exists k\in\mathbb{Z}\colon k^2 < 0\).}
{Es una proposición lógica falsa.}{Es una proposición lógica verdadera.}{No es una proposición lógica.}{Es imposible determinar si es una proposición lógica verdadera o falsa.}{}{}}
\End

\Data\ID{1103055701}
\Updated{2020-07-15 03:07:10}
\Level{B}
\SubArea{Logic and Sets}
\LANGen{
\Question{
Which diagram represents the interval that is the set of all solutions of the inequality \( -4 < x+1< 4 \).}
{}{}{}{}{}{}}
\LANGpl{
\Question{
Wskaż rysunek, na którym przedstawiono przedział, będący zbiorem wszystkich rozwiązań nierówności: \( -4 < x+1< 4 \).}
{}{}{}{}{}{}}
\LANGcs{
\Question{
Který obrázek znázorňuje množinu všech řešení nerovnice \( -4 < x+1< 4 \)?}
{}{}{}{}{}{}}
\LANGsk{
\Question{
Ktorý obrázok znázorňuje interval, kde je množina všetkých riešení nerovnice \( -4 < x+1< 4 \).}
{}{}{}{}{}{}}
\LANGes{
\Question{
Qué diagrama representa el intervalo que es el conjunto de todas las soluciones de la inecuación  \( -4 < x+1< 4 \).}
{}{}{}{}{}{}}

\IMGanswer{1103055701a.pdf}
\IMGanswer{1103055701b.pdf}
\IMGanswer{1103055701c.pdf}
\IMGanswer{1103055701d.pdf}\End

\Data\ID{2010000601}
\Updated{2021-12-01 04:12:21}
\Level{B}
\SubArea{Logic and Sets}
\LANGen{
\Question{
Given the sets \(A = \{x\in \mathbb{Z}:x> - 1\}\)
and \(B = \{x\in \mathbb{N}: x\leq 3\}\) find the intersection \(A \cap B\).}
{\(\{1;2;3\}\)}{\(\{0;1;2;3\}\)}{\(\{1;2\}\)}{\(\{-2;-1;0;1;2;3\}\)}{}{}}
\LANGpl{
\Question{
Dane są zbiory \(A = \{x\in \mathbb{Z}:x> - 1\}\)
i \(B = \{x\in \mathbb{N}: x\leq 3\}\).  Znajdź punkt przecięcia się \(A \cap B\).}
{\(\{1;2;3\}\)}{\(\{0;1;2;3\}\)}{\(\{1;2\}\)}{\(\{-2;-1;0;1;2;3\}\)}{}{}}
\LANGcs{
\Question{
Určete průnik množin \(A \cap B\), když \(A = \{x\in \mathbb{Z}:x> - 1\}\)
a \(B = \{x\in \mathbb{N}: x\leq 3\}\).}
{\(\{1;2;3\}\)}{\(\{0;1;2;3\}\)}{\(\{1;2\}\)}{\(\{-2;-1;0;1;2;3\}\)}{}{}}
\LANGsk{
\Question{
Určte prienik množín \(A \cap B\), ak \(A = \{x\in \mathbb{Z}:x> - 1\}\),  \(B = \{x\in \mathbb{N}: x\leq 3\}\).}
{\(\{1;2;3\}\)}{\(\{0;1;2;3\}\)}{\(\{1;2\}\)}{\(\{-2;-1;0;1;2;3\}\)}{}{}}
\LANGes{
\Question{
Dados los conjuntos \(A = \{x\in \mathbb{Z}:x> - 1\}\)
y \(B = \{x\in \mathbb{N}: x\leq 3\}\). Calcula la intersección \(A \cap B\).}
{\(\{1;2;3\}\)}{\(\{0;1;2;3\}\)}{\(\{1;2\}\)}{\(\{-2;-1;0;1;2;3\}\)}{}{}}
\End

\Data\ID{2010001401}
\Updated{2021-12-01 04:12:00}
\Level{B}
\SubArea{Logic and Sets}
\LANGen{
\Question{
Given the sets \(A =\mathbb{Z}\) 
and \(B = \{x\in \mathbb{N}\ \colon x < 5\}\),
find the union \(A\cup B\).}
{\(\mathbb{Z}\)}{\(\mathbb{N}\)}{\(\emptyset \)}{\( \{ 1;2;3;4\}\)}{}{}}
\LANGpl{
\Question{
Dane są \(A =\mathbb{Z}\) 
i \(B = \{x\in \mathbb{N}\ \colon x < 5\}\).
Wyznacz sumę zbiorów \(A\cup B\).}
{\(\mathbb{Z}\)}{\(\mathbb{N}\)}{\(\emptyset \)}{\( \{ 1;2;3;4\}\)}{}{}}
\LANGcs{
\Question{
Nechť \(A =\mathbb{Z}\) 
a \(B = \{x\in \mathbb{N}\ \colon x < 5\}\).
Určete sjednocení \(A\cup B\).}
{\(\mathbb{Z}\)}{\(\mathbb{N}\)}{\(\emptyset \)}{\( \{ 1;2;3;4\}\)}{}{}}
\LANGsk{
\Question{
Nech \(A =\mathbb{Z}\) 
a \(B = \{x\in \mathbb{N}\ \colon x < 5\}\),
potom určte zjednotenie \(A\cup B\).}
{\(\mathbb{Z}\)}{\(\mathbb{N}\)}{\(\emptyset \)}{\( \{ 1;2;3;4\}\)}{}{}}
\LANGes{
\Question{
Dados los conjuntos   \(A =\mathbb{Z}\) 
y \(B = \{x\in \mathbb{N}\ \colon x < 5\}\). Determina la unión \(A\cup B\).}
{\(\mathbb{Z}\)}{\(\mathbb{N}\)}{\(\emptyset \)}{\( \{ 1;2;3;4\}\)}{}{}}
\End

\Data\ID{2010001402}
\Updated{2021-12-01 04:12:06}
\Level{B}
\SubArea{Logic and Sets}
\LANGen{
\Question{
Find the set \( (B\setminus A) \cap C \) for \( A=[ -5;0) \), \( B=[ -1;10) \) and \( C=(-2;2 ] \).}
{\( [ 0;2]\)}{\( ( 0;2]\)}{\( [ -2;-1)\)}{\( [ -5;2]\)}{}{}}
\LANGpl{
\Question{
Wyznacz zbiór \( (B\setminus A) \cap C \), jeśli \( A=\langle -5;0) \), \( B=\langle -1;10) \) i \( C=(-2;2 \rangle \).}
{\( \langle 0;2\rangle\)}{\( ( 0;2\rangle\)}{\( \langle -2;-1)\)}{\( \langle -5;2\rangle\)}{}{}}
\LANGcs{
\Question{
Najděte množinu \( (B\setminus A) \cap C \), jestliže \( A=\langle -5;0) \), \( B=\langle -1;10) \) a \( C=(-2;2 \rangle \).}
{\( \langle 0;2\rangle\)}{\( ( 0;2\rangle\)}{\( \langle -2;-1)\)}{\( \langle -5;2\rangle\)}{}{}}
\LANGsk{
\Question{
Určte množinu \( (B\setminus A) \cap C \) pre \( A=\langle -5;0) \), \( B=\langle -1;10) \) a \( C=(-2;2 \rangle \).}
{\( \langle 0;2\rangle\)}{\( ( 0;2\rangle\)}{\( \langle -2;-1)\)}{\( \langle -5;2\rangle\)}{}{}}
\LANGes{
\Question{
Determina el conjunto \( (B\setminus A) \cap C \) para \( A=[ -5;0) \), \( B=[ -1;10) \) y \( C=(-2;2 ] \).}
{\( [ 0;2]\)}{\( ( 0;2]\)}{\( [ -2;-1)\)}{\( [ -5;2]\)}{}{}}
\End

\Data\ID{2010001405}
\Updated{2021-12-01 04:12:22}
\Level{B}
\SubArea{Logic and Sets}
\LANGen{
\Question{
Find the set difference \(A\setminus B\) 
for \(A = \{x\in \mathbb{Z}\ \colon \left |x\right | < 3\}\) 
and \(B = \{x\in \mathbb{N}\ \colon x \geq 2\}\).}
{\(\{ - 2;-1;0;1\}\)}{\( \emptyset \)}{\(\{ 2\}\)}{\(\{ - 2;-1;0\}\)}{}{}}
\LANGpl{
\Question{
Wyznacz róznicę zbiorów \(A\setminus B\) 
jeśli \(A = \{x\in \mathbb{Z}\ \colon \left |x\right | < 3\}\) 
i \(B = \{x\in \mathbb{N}\ \colon x \geq 2\}\).}
{\(\{ - 2;-1;0;1\}\)}{\( \emptyset \)}{\(\{ 2\}\)}{\(\{ - 2;-1;0\}\)}{}{}}
\LANGcs{
\Question{
Určete rozdíl \(A\setminus B\), jestliže \(A = \{x\in \mathbb{Z}\ \colon \left |x\right | < 3\}\) 
a \(B = \{x\in \mathbb{N}\ \colon x \geq 2\}\).}
{\(\{ - 2;-1;0;1\}\)}{\( \emptyset \)}{\(\{ 2\}\)}{\(\{ - 2;-1;0\}\)}{}{}}
\LANGsk{
\Question{
Určte rozdiel množín \(A\setminus B\) 
pre \(A = \{x\in \mathbb{Z}\ \colon \left |x\right | < 3\}\) 
a \(B = \{x\in \mathbb{N}\ \colon x \geq 2\}\).}
{\(\{ - 2;-1;0;1\}\)}{\( \emptyset \)}{\(\{ 2\}\)}{\(\{ - 2;-1;0\}\)}{}{}}
\LANGes{
\Question{
Determina la diferencia de conjuntos \(A\setminus B\) 
para \(A = \{x\in \mathbb{Z}\ \colon \left |x\right | < 3\}\) 
y \(B = \{x\in \mathbb{N}\ \colon x \geq 2\}\).}
{\(\{ - 2;-1;0;1\}\)}{\( \emptyset \)}{\(\{ 2\}\)}{\(\{ - 2;-1;0\}\)}{}{}}
\End

\Data\ID{9000080901}
\Updated{2020-07-15 03:07:37}
\Level{B}
\SubArea{Logic and Sets}
\LANGen{
\Question{
Find the intersection \(A \cap B\) 
for \(A = \{ - 5;0;1.5;2;6\}\) 
and \(B = \{x\in \mathbb{Z};x\geq 0\}\).}
{\(\{0;2;6\}\)}{\(\{0;1.5;2;6\}\)}{\(\{1.5;2;6\}\)}{\(\mathbb{Z}\)}{}{}}
\LANGpl{
\Question{
Wyznacz iloczyn zbiorów \(A \cap B\), gdzie \(A = \{ - 5;0;1.5;2;6\}\) i \(B = \{x\in \mathbb{Z};x\geq 0\}\).}
{\(\{0;2;6\}\)}{\(\{0;1.5;2;6\}\)}{\(\{1.5;2;6\}\)}{\(\mathbb{Z}\)}{}{}}
\LANGcs{
\Question{
Určete průnik množin \(A\),
\(B\), jestliže
\(A = \{ - 5;0;1{,}5;2;6\}\),
\(B = \{x\in \mathbb{Z};x\geq 0\}\).}
{\(\{0;2;6\}\)}{\(\{0;1{,}5;2;6\}\)}{\(\{1{,}5;2;6\}\)}{\(\mathbb{Z}\)}{}{}}
\LANGsk{
\Question{
Určte prienik množín \(A\),
\(B\), ak
\(A = \{ - 5;0;1{,}5;2;6\}\),
\(B = \{x\in \mathbb{Z};x\geq 0\}\).}
{\(\{0;2;6\}\)}{\(\{0;1{,}5;2;6\}\)}{\(\{1{,}5;2;6\}\)}{\(\mathbb{Z}\)}{}{}}
\LANGes{
\Question{
Calcula la intersección  \(A \cap B\) 
para \(A = \{ - 5;0;1.5;2;6\}\) 
y \(B = \{x\in \mathbb{Z};x\geq 0\}\).}
{\(\{0;2;6\}\)}{\(\{0;1.5;2;6\}\)}{\(\{1.5;2;6\}\)}{\(\mathbb{Z}\)}{}{}}
\End

\Data\ID{9000080902}
\Updated{2022-04-23 05:04:02}
\Level{B}
\SubArea{Logic and Sets}
\LANGen{
\Question{
Given the sets \(A = \{x\in \mathbb{Z}:x\geq - 2\}\)
and \(B = \{x\in \mathbb{N}: x\leq 5\}\) find the intersection \(A \cap B\).}
{\(\{1;2;3;4;5\}\)}{\(\{0;1;2;3;4;5\}\)}{\(\{0;1;2;3;4\}\)}{\(\{ - 2;-1;0;1;2;3;4;5\}\)}{}{}}
\LANGpl{
\Question{
Wyznacz iloczyn zbiorów \(A \cap B\), gdzie  \(A = \{x\in \mathbb{Z}:x\geq - 2\}\) i \(B = \{x\in \mathbb{N}:x\leq 5\}\).}
{\(\{1;2;3;4;5\}\)}{\(\{0;1;2;3;4;5\}\)}{\(\{0;1;2;3;4\}\)}{\(\{ - 2;-1;0;1;2;3;4;5\}\)}{}{}}
\LANGcs{
\Question{
Určete průnik množin \(A\),
\(B\), jestliže
\(A = \{x\in \mathbb{Z}:x\geq - 2\}\),
\(B = \{x\in \mathbb{N}:x\leq 5\}\).}
{\(\{1;2;3;4;5\}\)}{\(\{0;1;2;3;4;5\}\)}{\(\{0;1;2;3;4\}\)}{\(\{ - 2;-1;0;1;2;3;4;5\}\)}{}{}}
\LANGsk{
\Question{
Určte prienik množín \(A\),
\(B\), ak
\(A = \{x\in \mathbb{Z}:x\geq - 2\}\),
\(B = \{x\in \mathbb{N}:x\leq 5\}\).}
{\(\{1;2;3;4;5\}\)}{\(\{0;1;2;3;4;5\}\)}{\(\{0;1;2;3;4\}\)}{\(\{ - 2;-1;0;1;2;3;4;5\}\)}{}{}}
\LANGes{
\Question{
Calcula la intersección \(A \cap B\) 
para \(A = \{x\in \mathbb{Z}:x\geq - 2\}\)
y \(B = \{x\in \mathbb{N}:x\leq 5\}\).}
{\(\{1;2;3;4;5\}\)}{\(\{0;1;2;3;4;5\}\)}{\(\{0;1;2;3;4\}\)}{\(\{ - 2;-1;0;1;2;3;4;5\}\)}{}{}}
\End

\Data\ID{9000080903}
\Updated{2020-07-15 03:07:37}
\Level{B}
\SubArea{Logic and Sets}
\LANGen{
\Question{
Find the union \(A\cup B\) 
for \(A = \{x\in \mathbb{Z};x\geq - 3\}\) 
and \(B = \{x\in \mathbb{N};x < 8\}\).}
{\(\{x\in \mathbb{Z};x\geq - 3\}\)}{\(\{1;2;3;4;5;6;7\}\)}{\(\{ - 3;-2;-1;0;1;2;3;4;5;6;7\}\)}{\(\mathbb{Z}\)}{}{}}
\LANGpl{
\Question{
Wyznacz sumę zbiorów  \(A\cup B\), jeśli \(A = \{x\in \mathbb{Z};x\geq - 3\}\) i \(B = \{x\in \mathbb{N};x < 8\}\).}
{\(\{x\in \mathbb{Z};x\geq - 3\}\)}{\(\{1;2;3;4;5;6;7\}\)}{\(\{ - 3;-2;-1;0;1;2;3;4;5;6;7\}\)}{\(\mathbb{Z}\)}{}{}}
\LANGcs{
\Question{
Určete sjednocení množin \(A\),
\(B\), jestliže
\(A = \{x\in \mathbb{Z};x\geq - 3\}\),
\(B = \{x\in \mathbb{N};x < 8\}\).}
{\(\{x\in \mathbb{Z};x\geq - 3\}\)}{\(\{1;2;3;4;5;6;7\}\)}{\(\{ - 3;-2;-1;0;1;2;3;4;5;6;7\}\)}{\(\mathbb{Z}\)}{}{}}
\LANGsk{
\Question{
Určte zjednotenie množín \(A\),
\(B\), ak
\(A = \{x\in \mathbb{Z};x\geq - 3\}\),
\(B = \{x\in \mathbb{N};x < 8\}\).}
{\(\{x\in \mathbb{Z};x\geq - 3\}\)}{\(\{1;2;3;4;5;6;7\}\)}{\(\{ - 3;-2;-1;0;1;2;3;4;5;6;7\}\)}{\(\mathbb{Z}\)}{}{}}
\LANGes{
\Question{
Calcula la unión \(A\cup B\) 
para \(A = \{x\in \mathbb{Z};x\geq - 3\}\) 
y \(B = \{x\in \mathbb{N};x < 8\}\).}
{\(\{x\in \mathbb{Z};x\geq - 3\}\)}{\(\{1;2;3;4;5;6;7\}\)}{\(\{ - 3;-2;-1;0;1;2;3;4;5;6;7\}\)}{\(\mathbb{Z}\)}{}{}}
\End

\Data\ID{9000080904}
\Updated{2022-04-23 05:04:08}
\Level{B}
\SubArea{Logic and Sets}
\LANGen{
\Question{
Given the sets \(A =\mathbb{N}\) 
and \(B = \{x\in \mathbb{Z}\colon x > 8\}\),
find the union \(A\cup B\).}
{\(\mathbb{N}\)}{\(\emptyset \)}{\(\{x\in \mathbb{Z};x > 8\}\)}{\(\mathbb{Z}\)}{}{}}
\LANGpl{
\Question{
Wyznacz sumę zbiorów  \(A\cup B\), jeśli  \(A =\mathbb{N}\) i \(B = \{x\in \mathbb{Z};x > 8\}\).}
{\(\mathbb{N}\)}{\(\emptyset \)}{\(\{x\in \mathbb{Z};x > 8\}\)}{\(\mathbb{Z}\)}{}{}}
\LANGcs{
\Question{
Určete sjednocení množin \(A\),
\(B\), jestliže
\(A =\mathbb{N}\),
\(B = \{x\in \mathbb{Z};x > 8\}\).}
{\(\mathbb{N}\)}{\(\emptyset \)}{\(\{x\in \mathbb{Z};x > 8\}\)}{\(\mathbb{Z}\)}{}{}}
\LANGsk{
\Question{
Určte zjednotenie množín \(A\),
\(B\), ak
\(A =\mathbb{N}\),
\(B = \{x\in \mathbb{Z};x > 8\}\).}
{\(\mathbb{N}\)}{\(\emptyset \)}{\(\{x\in \mathbb{Z};x > 8\}\)}{\(\mathbb{Z}\)}{}{}}
\LANGes{
\Question{
Calcula la unión  \(A\cup B\) 
para \(A =\mathbb{N}\) 
y \(B = \{x\in \mathbb{Z};x > 8\}\).}
{\(\mathbb{N}\)}{\(\emptyset \)}{\(\{x\in \mathbb{Z};x > 8\}\)}{\(\mathbb{Z}\)}{}{}}
\End

\Data\ID{9000080905}
\Updated{2020-07-15 03:07:37}
\Level{B}
\SubArea{Logic and Sets}
\LANGen{
\Question{
Identify the set \(B\) 
which satisfies 
\[ 
 A\cup B = C
\]
if \(A = \{x\in \mathbb{N};x < 3\}\) and
\(C = \{0;1;2\}\).}
{\(\{0;1;2\},\ \{0;1\},\ \{0;2\},\ \{0\}\)}{no solution exist}{\(\emptyset \)}{\(\{0;1;2\},\ \{0;1\},\ \{1;2\},\ \{0;2\}\)}{}{}}
\LANGpl{
\Question{
Wyznacz zbiór \(B\), dla którego zachodzi równość \[ A\cup B = C\text{,} \]
jeśli  \(A = \{x\in \mathbb{N};x < 3\}\) i
\(C = \{0;1;2\}\).}
{\(\{0;1;2\},\ \{0;1\},\ \{0;2\},\ \{0\}\)}{Rozwiązanie nie istnieje.}{\(\emptyset \)}{\(\{0;1;2\},\ \{0;1\},\ \{1;2\},\ \{0;2\}\)}{}{}}
\LANGcs{
\Question{
Určete všechny množiny \(B\),
pro které platí: \(A\cup B = C\),
jestliže \(A = \{x\in \mathbb{N};x < 3\}\) 
a \(C = \{0;1;2\}\).}
{\(\{0;1;2\},\ \{0;1\},\ \{0;2\},\ \{0\}\)}{řešení neexistuje}{\(\emptyset \)}{\(\{0;1;2\},\ \{0;1\},\ \{1;2\},\ \{0;2\}\)}{}{}}
\LANGsk{
\Question{
Určte všetky množiny \(B\),
pre ktoré platí: \(A\cup B = C\),
ak \(A = \{x\in \mathbb{N};x < 3\}\) 
a \(C = \{0;1;2\}\).}
{\(\{0;1;2\},\ \{0;1\},\ \{0;2\},\ \{0\}\)}{riešenie neexistuje}{\(\emptyset \)}{\(\{0;1;2\},\ \{0;1\},\ \{1;2\},\ \{0;2\}\)}{}{}}
\LANGes{
\Question{
Identifica el conjunto \(B\) 
que satisface  
\[ 
 A\cup B = C
\]
si \(A = \{x\in \mathbb{N};x < 3\}\) y
\(C = \{0;1;2\}\).}
{\(\{0;1;2\},\ \{0;1\},\ \{0;2\},\ \{0\}\)}{No existe ninguna solución.}{\(\emptyset \)}{\(\{0;1;2\},\ \{0;1\},\ \{1;2\},\ \{0;2\}\)}{}{}}
\End

\Data\ID{9000080906}
\Updated{2020-07-15 03:07:37}
\Level{B}
\SubArea{Logic and Sets}
\LANGen{
\Question{
Find \(B'_{A}\) (the
complement to \(B\) 
in \(A\))
for \(A = \{x\in \mathbb{N};x < 9\}\) 
and \(B = \{4;5;6;7\}\).}
{\(\{1;2;3;8\}\)}{\(\emptyset \)}{\(\{4;5;6;7\}\)}{\(\{0;1;2;3;8\}\)}{}{}}
\LANGpl{
\Question{
Wyznacz zbiór  \(B'_{A}\) (dopełnienie zbioru \(B\) 
w zbiorze \(A\)), jeśli \(A = \{x\in \mathbb{N};x < 9\}\) 
i \(B = \{4;5;6;7\}\).}
{\(\{1;2;3;8\}\)}{\(\emptyset \)}{\(\{4;5;6;7\}\)}{\(\{0;1;2;3;8\}\)}{}{}}
\LANGcs{
\Question{
Určete množinu \(B'_{A}\) 
(doplněk množiny \(B\) 
v množině \(A\)),
jestliže \(A = \{x\in \mathbb{N};x < 9\}\),
\(B = \{4;5;6;7\}\).}
{\(\{1;2;3;8\}\)}{\(\emptyset \)}{\(\{4;5;6;7\}\)}{\(\{0;1;2;3;8\}\)}{}{}}
\LANGsk{
\Question{
Určte množinu \(B'_{A}\) 
(doplnok množiny \(B\) 
v množine \(A\)),
ak \(A = \{x\in \mathbb{N};x < 9\}\),
\(B = \{4;5;6;7\}\).}
{\(\{1;2;3;8\}\)}{\(\emptyset \)}{\(\{4;5;6;7\}\)}{\(\{0;1;2;3;8\}\)}{}{}}
\LANGes{
\Question{
Calcula \(B'_{A}\) (el complementario del conjunto \(B\) 
en el conjunto \(A\))
para \(A = \{x\in \mathbb{N};x < 9\}\) 
y \(B = \{4;5;6;7\}\).}
{\(\{1;2;3;8\}\)}{\(\emptyset \)}{\(\{4;5;6;7\}\)}{\(\{0;1;2;3;8\}\)}{}{}}
\End

\Data\ID{9000080907}
\Updated{2020-07-15 03:07:37}
\Level{B}
\SubArea{Logic and Sets}
\LANGen{
\Question{
Find \(B'_{A}\) (the
complement to \(B\) 
in \(A\))
for \(A =\mathbb{Z}\) 
and \(B = \{x\in \mathbb{Z};\left |x\right | > 3\}\).}
{\(\{ - 3;-2;-1;0;1;2;3\}\)}{\(\{ - 2;-1;0;1;2\}\)}{\(\{0;1;2;3\}\)}{\(\{1;2;3\}\)}{}{}}
\LANGpl{
\Question{
Wyznacz zbiór \(B'_{A}\) (dopełnienie zbioru \(B\) 
w zbiorze  \(A\)), jeśli \(A =\mathbb{Z}\) 
i \(B = \{x\in \mathbb{Z};\left |x\right | > 3\}\).}
{\(\{ - 3;-2;-1;0;1;2;3\}\)}{\(\{ - 2;-1;0;1;2\}\)}{\(\{0;1;2;3\}\)}{\(\{1;2;3\}\)}{}{}}
\LANGcs{
\Question{
Určete množinu \(B'_{A}\) 
(doplněk množiny \(B\) 
v množině \(A\)),
jestliže \(A =\mathbb{Z}\),
\(B = \{x\in \mathbb{Z};\left |x\right | > 3\}\).}
{\(\{ - 3;-2;-1;0;1;2;3\}\)}{\(\{ - 2;-1;0;1;2\}\)}{\(\{0;1;2;3\}\)}{\(\{1;2;3\}\)}{}{}}
\LANGsk{
\Question{
Určte množinu \(B'_{A}\) 
(doplnok množiny \(B\) 
v množine \(A\)),
ak \(A =\mathbb{Z}\),
\(B = \{x\in \mathbb{Z};\left |x\right | > 3\}\).}
{\(\{ - 3;-2;-1;0;1;2;3\}\)}{\(\{ - 2;-1;0;1;2\}\)}{\(\{0;1;2;3\}\)}{\(\{1;2;3\}\)}{}{}}
\LANGes{
\Question{
Calcula \(B'_{A}\) (el complementario del conjunto  \(B\) 
en el conjunto \(A\))
para \(A =\mathbb{Z}\) 
y \(B = \{x\in \mathbb{Z};\left |x\right | > 3\}\).}
{\(\{ - 3;-2;-1;0;1;2;3\}\)}{\(\{ - 2;-1;0;1;2\}\)}{\(\{0;1;2;3\}\)}{\(\{1;2;3\}\)}{}{}}
\End

\Data\ID{9000080908}
\Updated{2020-07-15 03:07:37}
\Level{B}
\SubArea{Logic and Sets}
\LANGen{
\Question{
Find the set difference \(A\setminus B\) 
for \(A = \{ - 2;-1;0;1;2\}\) 
and \(B = \{x\in \mathbb{Z};x < 2\}\).}
{\(\{2\}\)}{\(\{ - 2;-1;0;1;2\}\)}{\(\{0;1\}\)}{\(\emptyset \)}{}{}}
\LANGpl{
\Question{
Wyznacz różnicę zbiorów \(A\setminus B\), jeśli \(A = \{ - 2;-1;0;1;2\}\) 
i \(B = \{x\in \mathbb{Z};x < 2\}\).}
{\(\{2\}\)}{\(\{ - 2;-1;0;1;2\}\)}{\(\{0;1\}\)}{\(\emptyset \)}{}{}}
\LANGcs{
\Question{
Určete rozdíl \(A\setminus B\),
jestliže \(A = \{ - 2;-1;0;1;2\}\),
\(B = \{x\in \mathbb{Z};x < 2\}\).}
{\(\{2\}\)}{\(\{ - 2;-1;0;1;2\}\)}{\(\{0;1\}\)}{\(\emptyset \)}{}{}}
\LANGsk{
\Question{
Určte rozdiel \(A\setminus B\),
ak \(A = \{ - 2;-1;0;1;2\}\),
\(B = \{x\in \mathbb{Z};x < 2\}\).}
{\(\{2\}\)}{\(\{ - 2;-1;0;1;2\}\)}{\(\{0;1\}\)}{\(\emptyset \)}{}{}}
\LANGes{
\Question{
Calcula la diferencia \(A\setminus B\) 
para \(A = \{ - 2;-1;0;1;2\}\) 
y \(B = \{x\in \mathbb{Z};x < 2\}\).}
{\(\{2\}\)}{\(\{ - 2;-1;0;1;2\}\)}{\(\{0;1\}\)}{\(\emptyset \)}{}{}}
\End

\Data\ID{9000080909}
\Updated{2020-07-15 03:07:37}
\Level{B}
\SubArea{Logic and Sets}
\LANGen{
\Question{
Find the set difference \(B\setminus A\) 
for \(A = \{x\in \mathbb{Z};x < 2\}\) 
and \(B = \{x\in \mathbb{Z};x < 5\}\).}
{\(\{2;3;4\}\)}{\(\{x\in \mathbb{Z};x < 2\}\)}{\(\{3;4\}\)}{\(\emptyset \)}{}{}}
\LANGpl{
\Question{
Wyznacz różnicę zbiorów \(B\setminus A\), jeśli \(A = \{x\in \mathbb{Z};x < 2\}\) 
i \(B = \{x\in \mathbb{Z};x < 5\}\).}
{\(\{2;3;4\}\)}{\(\{x\in \mathbb{Z};x < 2\}\)}{\(\{3;4\}\)}{\(\emptyset \)}{}{}}
\LANGcs{
\Question{
Určete rozdíl \(B\setminus A\),
jestliže \(A = \{x\in \mathbb{Z};x < 2\}\),
\(B = \{x\in \mathbb{Z};x < 5\}\).}
{\(\{2;3;4\}\)}{\(\{x\in \mathbb{Z};x < 2\}\)}{\(\{3;4\}\)}{\(\emptyset \)}{}{}}
\LANGsk{
\Question{
Určte rozdiel \(B\setminus A\),
ak \(A = \{x\in \mathbb{Z};x < 2\}\),
\(B = \{x\in \mathbb{Z};x < 5\}\).}
{\(\{2;3;4\}\)}{\(\{x\in \mathbb{Z};x < 2\}\)}{\(\{3;4\}\)}{\(\emptyset \)}{}{}}
\LANGes{
\Question{
Calcula la diferencia \(B\setminus A\) 
para \(A = \{x\in \mathbb{Z};x < 2\}\) 
y \(B = \{x\in \mathbb{Z};x < 5\}\).}
{\(\{2;3;4\}\)}{\(\{x\in \mathbb{Z};x < 2\}\)}{\(\{3;4\}\)}{\(\emptyset \)}{}{}}
\End

\Data\ID{9000080910}
\Updated{2022-04-23 05:04:08}
\Level{B}
\SubArea{Logic and Sets}
\LANGen{
\Question{
Find the set difference \(A\setminus B\) 
for \(A = \{x\in \mathbb{Z};\left |x\right | < 3\}\) 
and \(B = \{x\in \mathbb{N};x < 5\}\).}
{\(\{ - 2;-1;0\}\)}{\(\emptyset \)}{\(\{- 2;-1\}\)}{\(\{3;4\}\)}{}{}}
\LANGpl{
\Question{
Wyznacz różnicę zbiorów  \(A\setminus B\), jeśli \(A = \{x\in \mathbb{Z};\left |x\right | < 3\}\) 
i \(B = \{x\in \mathbb{N};x < 5\}\).}
{\(\{ - 2;-1;0\}\)}{\(\emptyset \)}{\(\{- 2;-1\}\)}{\(\{3;4\}\)}{}{}}
\LANGcs{
\Question{
Určete rozdíl \(A\setminus B\),
jestliže \(A = \{x\in \mathbb{Z};\left |x\right | < 3\}\),
\(B = \{x\in \mathbb{N};x < 5\}\).}
{\(\{ - 2;-1;0\}\)}{\(\emptyset \)}{\(\{- 2;-1\}\)}{\(\{3;4\}\)}{}{}}
\LANGsk{
\Question{
Určte rozdiel \(A\setminus B\),
ak \(A = \{x\in \mathbb{Z};\left |x\right | < 3\}\),
\(B = \{x\in \mathbb{N};x < 5\}\).}
{\(\{ - 2;-1;0\}\)}{\(\emptyset \)}{\(\{- 2;-1\}\)}{\(\{3;4\}\)}{}{}}
\LANGes{
\Question{
Calcula la diferencia  \(A\setminus B\) 
para \(A = \{x\in \mathbb{Z};\left |x\right | < 3\}\) 
y \(B = \{x\in \mathbb{N};x < 5\}\).}
{\(\{ - 2;-1;0\}\)}{\(\emptyset \)}{\(\{- 2;-1\}\)}{\(\{3;4\}\)}{}{}}
\End

\Data\ID{9000086601}
\Updated{2020-07-15 03:07:37}
\Level{B}
\SubArea{Logic and Sets}
\LANGen{
\Question{
Determine the truth values of propositions \(a\) and \(b\) if you know that the compound proposition
\[ 
 \neg (a \vee b)
\]
is true.}
{Both statements are false.}{Both statements are true.}{The statement \(a\) 
is true, \(b\) 
is false.}{The statement \(a\) 
is false, \(b\) 
is true.}{}{}}
\LANGpl{
\Question{
Oceń wartość logiczną zdań  \(a\) i \(b\), jeśli wiadomo, że zdanie
\[ 
 \neg (a \vee b)
\]
jest prawdziwe.}
{Oba zdania są fałszywe.}{Oba zdania są prawdziwe.}{Zdanie \(a\) jest prawdziwe, zdanie \(b\) jest fałszywe.}{Zdanie \(a\) jest fałszywe, zdanie \(b\) jest prawdziwe.}{}{}}
\LANGcs{
\Question{
Určete pravdivostní hodnoty výroků
\(a\) a
\(b\),
víte-li, že pravdivostní hodnota složeného výroku
\(\neg (a \vee b)\) je \(1\).}
{Pravdivostní hodnota \(a\) 
je \(0\), pravdivostní
hodnota \(b\) 
je \(0\).}{Pravdivostní hodnota \(a\) 
je \(1\), pravdivostní
hodnota \(b\) 
je \(1\).}{Pravdivostní hodnota \(a\) 
je \(1\), pravdivostní
hodnota \(b\) 
je \(0\).}{Pravdivostní hodnota \(a\) 
je \(0\), pravdivostní
hodnota \(b\) 
je \(1\).}{}{}}
\LANGsk{
\Question{
Určte pravdivostné hodnoty výrokov
\(a\) a
\(b\),
ak viete, že pravdivostná hodnota zloženého výroku
\(\neg (a \vee b)\) je \(1\).}
{Pravdivostná hodnota \(a\) 
je \(0\), pravdivostná
hodnota \(b\) 
je \(0\).}{Pravdivostná hodnota \(a\) 
je \(1\), pravdivostná
hodnota \(b\) 
je \(1\).}{Pravdivostná hodnota \(a\) 
je \(1\), pravdivostná
hodnota \(b\) 
je \(0\).}{Pravdivostná hodnota \(a\) 
je \(0\), pravdivostná
hodnota \(b\) 
je \(1\).}{}{}}
\LANGes{
\Question{
Determina los valores de verdad de las proposiciones \(a\) y \(b\) b sabiendo que la proposición compuesta
\[ 
 \neg (a \vee b)
\]
es verdadera.}
{Ambas proposiciones son falsas.}{Ambas proposiciones son verdaderas.}{La proposición  \(a\) 
es verdadera, \(b\) 
es falsa.}{La proposición  \(a\) 
es falsa, \(b\) 
es verdadera.}{}{}}
\End

\Data\ID{9000086602}
\Updated{2020-07-15 03:07:37}
\Level{B}
\SubArea{Logic and Sets}
\LANGen{
\Question{
Determine the truth values of propositions \(a\) and \(b\) if you know that the compound proposition
\[ 
 \neg a \vee b
\]
is false.}
{The statement \(a\) 
is true, \(b\) 
is false.}{Both statements are true.}{The statement \(a\) 
is false, \(b\) 
is true.}{Both statements are false.}{}{}}
\LANGpl{
\Question{
Oceń wartość logiczną zdań \(a\) 
i \(b\), 
jeśli wiadomo, że zdanie
\[ 
 \neg a \vee b
\]
jest fałszywe.}
{Zdanie \(a\) jest prawdziwe, zdanie \(b\) jest fałszywe.}{Oba zdania są prawdziwe.}{Zdanie \(a\) jest fałszywe, zdanie \(b\) jest prawdziwe.}{Oba zdania są fałszywe.}{}{}}
\LANGcs{
\Question{
Určete pravdivostní hodnoty výroků
\(a\) a
\(b\),
víte-li, že pravdivostní hodnota složeného výroku
\(\neg a \vee b\) je
\(0\).}
{Pravdivostní hodnota \(a\) 
je \(1\), pravdivostní
hodnota \(b\) 
je \(0\).}{Pravdivostní hodnota \(a\) 
je \(1\), pravdivostní
hodnota \(b\) 
je \(1\).}{Pravdivostní hodnota \(a\) 
je \(0\), pravdivostní
hodnota \(b\) 
je \(1\).}{Pravdivostní hodnota \(a\) 
je \(0\), pravdivostní
hodnota \(b\) 
je \(0\).}{}{}}
\LANGsk{
\Question{
Určte pravdivostné hodnoty výrokov
\(a\) a
\(b\),
ak viete, že pravdivostná hodnota zloženého výroku
\(\neg a \vee b\) je
\(0\).}
{Pravdivostná hodnota \(a\) 
je \(1\), pravdivostná
hodnota \(b\) 
je \(0\).}{Pravdivostná hodnota \(a\) 
je \(1\), pravdivostná
hodnota \(b\) 
je \(1\).}{Pravdivostná hodnota \(a\) 
je \(0\), pravdivostná
hodnota \(b\) 
je \(1\).}{Pravdivostná hodnota \(a\) 
je \(0\), pravdivostná
hodnota \(b\) 
je \(0\).}{}{}}
\LANGes{
\Question{
Determina los valores de verdad de las proposiciones \(a\) y \(b\) sabiendo que la proposición compuesta 
\[ 
 \neg a \vee b
\]
es falsa.}
{La proposición \(a\) 
es verdadera, \(b\) 
es falsa.}{Ambas proposiciones son verdaderas.}{La proposición \(a\) 
es falsa, \(b\) 
es verdadera.}{Ambas proposiciones son falsas.}{}{}}
\End

\Data\ID{9000086604}
\Updated{2020-07-15 03:07:37}
\Level{B}
\SubArea{Logic and Sets}
\LANGen{
\Question{
Determine the truth values of propositions \(a\) and \(b\) if you know that the compound proposition
\[ 
 \neg (a \wedge \neg b)
\]
is false.}
{The statement \(a\) 
is true, \(b\) 
is false.}{Both statements are true.}{The statement \(a\) 
is false, \(b\) 
is true.}{Both statements are false.}{}{}}
\LANGpl{
\Question{
Oceń wartość logiczną zdań \(a\) i \(b\), 
jeśli wiadomo, że zdanie
\[ 
 \neg (a \wedge \neg b)
\]
jest fałszywe.}
{Zdanie \(a\) jest prawdziwe, zdanie \(b\) jest fałszywe.}{Oba zdania są prawdziwe.}{Zdanie \(a\) jest fałszywe, zdanie \(b\) jest prawdziwe.}{Oba zdania są fałszywe.}{}{}}
\LANGcs{
\Question{
Určete pravdivostní hodnoty výroků
\(a\) a
\(b\),
víte-li, že pravdivostní hodnota složeného výroku
\(\neg (a \wedge \neg b)\) je
\(0\).}
{Pravdivostní hodnota \(a\) 
je \(1\), pravdivostní
hodnota \(b\) 
je \(0\).}{Pravdivostní hodnota \(a\) 
je \(1\), pravdivostní
hodnota \(b\) 
je \(1\).}{Pravdivostní hodnota \(a\) 
je \(0\), pravdivostní
hodnota \(b\) 
je \(1\).}{Pravdivostní hodnota \(a\) 
je \(0\), pravdivostní
hodnota \(b\) 
je \(0\).}{}{}}
\LANGsk{
\Question{
Určte pravdivostné hodnoty výrokov
\(a\) a
\(b\),
ak viete, že pravdivostná hodnota zloženého výroku
\(\neg (a \wedge \neg b)\) je
\(0\).}
{Pravdivostná hodnota \(a\) 
je \(1\), pravdivostná
hodnota \(b\) 
je \(0\).}{Pravdivostná hodnota \(a\) 
je \(1\), pravdivostná
hodnota \(b\) 
je \(1\).}{Pravdivostná hodnota \(a\) 
je \(0\), pravdivostná
hodnota \(b\) 
je \(1\).}{Pravdivostná hodnota \(a\) 
je \(0\), pravdivostná
hodnota \(b\) 
je \(0\).}{}{}}
\LANGes{
\Question{
Determina los valores de verdad de las proposiciones  \(a\) y \(b\) sabiendo que la proposición compuesta
\[ 
 \neg (a \wedge \neg b)
\]
es falsa.}
{La proposición \(a\) 
es verdadera, \(b\) 
es falsa.}{Ambas proposiciones son verdaderas.}{La proposición  \(a\) 
es falsa, \(b\) 
es verdadera.}{Ambas proposiciones son falsas.}{}{}}
\End

\Data\ID{9000086605}
\Updated{2020-07-15 03:07:37}
\Level{B}
\SubArea{Logic and Sets}
\LANGen{
\Question{
Determine the truth values of propositions \(a\) and \(b\) if you know that the compound proposition 
\[ 
 \neg a\implies \neg b
\]
is false.}
{The statement \(a\) 
is false, \(b\) 
is true.}{Both statements are true.}{The statement \(a\) 
is true, \(b\) 
is false.}{Both statements are false.}{}{}}
\LANGpl{
\Question{
Oceń wartość logiczną zdań  \(a\) 
i  \(b\),
jeśli wiadomo, że zdanie 
\[ 
 \neg a\implies \neg b
\]
jest fałszywe.}
{Zdanie \(a\) jest fałszywe, zdanie \(b\) jest prawdziwe.}{Oba zdania są prawdziwe.}{Zdanie \(a\) jest prawdziwe, zdanie \(b\) jest fałszywe.}{Oba zdania są fałszywe.}{}{}}
\LANGcs{
\Question{
Určete pravdivostní hodnoty výroků
\(a\) a
\(b\),
víte-li, že pravdivostní hodnota složeného výroku
\(\neg a\implies \neg b\) je
\(0\).}
{Pravdivostní hodnota \(a\) 
je \(0\), pravdivostní
hodnota \(b\) 
je \(1\).}{Pravdivostní hodnota \(a\) 
je \(1\), pravdivostní
hodnota \(b\) 
je \(1\).}{Pravdivostní hodnota \(a\) 
je \(1\), pravdivostní
hodnota \(b\) 
je \(0\).}{Pravdivostní hodnota \(a\) 
je \(0\), pravdivostní
hodnota \(b\) 
je \(0\).}{}{}}
\LANGsk{
\Question{
Určte pravdivostné hodnoty výrokov
\(a\) a
\(b\),
ak viete, že pravdivostná hodnota zloženého výroku
\(\neg a\implies \neg b\) je
\(0\).}
{Pravdivostná hodnota \(a\) 
je \(0\), pravdivostná
hodnota \(b\) 
je \(1\).}{Pravdivostná hodnota \(a\) 
je \(1\), pravdivostná
hodnota \(b\) 
je \(1\).}{Pravdivostná hodnota \(a\) 
je \(1\), pravdivostná
hodnota \(b\) 
je \(0\).}{Pravdivostná hodnota \(a\) 
je \(0\), pravdivostná
hodnota \(b\) 
je \(0\).}{}{}}
\LANGes{
\Question{
Determina los valores de verdad de las proposiciones  \(a\) y \(b\) sabiendo que la proposición compuesta 
\[ 
 \neg a\implies \neg b
\]
es falsa.}
{La proposición  \(a\) 
es falsa, \(b\) 
es verdadera.}{Ambas proposiciones son verdaderas.}{La proposición  \(a\) 
es verdadera, \(b\) 
es falsa.}{Ambas proposiciones son falsas.}{}{}}
\End

\Data\ID{9000086606}
\Updated{2022-04-23 05:04:57}
\Level{B}
\SubArea{Logic and Sets}
\LANGen{
\Question{
Determine the truth values of statements \(a\) and \(b\) if you know that the compound statement
\[ 
 a \iff (a \vee b)
\]
is false.}
{The statement \(a\) 
is false, \(b\) 
is true.}{Both statements are true.}{The statement \(a\) 
is true, \(b\) 
is false.}{Both statements are false.}{}{}}
\LANGpl{
\Question{
Oceń wartość logiczną zdań \(a\) 
i \(b\), jeśli wiadomo, że zdanie  
\[ 
 a \iff (a \vee b)
\]
jest fałszywe.}
{Zdanie \(a\) jest fałszywe, zdanie \(b\) jest prawdziwe.}{Oba zdania są prawdziwe.}{Zdanie \(a\) jest prawdziwe, zdanie \(b\) jest fałszywe.}{Oba zdania są fałszywe.}{}{}}
\LANGcs{
\Question{
Určete pravdivostní hodnoty výroků
\(a\) a
\(b\),
víte-li, že pravdivostní hodnota složeného výroku
\(a \iff (a \vee b)\) je
\(0\).}
{Pravdivostní hodnota \(a\) 
je \(0\), pravdivostní
hodnota \(b\) 
je \(1\).}{Pravdivostní hodnota \(a\) 
je \(1\), pravdivostní
hodnota \(b\) 
je \(1\).}{Pravdivostní hodnota \(a\) 
je \(1\), pravdivostní
hodnota \(b\) 
je \(0\).}{Pravdivostní hodnota \(a\) 
je \(0\), pravdivostní
hodnota \(b\) 
je \(0\).}{}{}}
\LANGsk{
\Question{
Určte pravdivostné hodnoty výrokov
\(a\) a
\(b\),
ak viete, že pravdivostná hodnota zloženého výroku
\(a \iff (a \vee b)\) je
\(0\).}
{Pravdivostná hodnota \(a\) 
je \(0\), pravdivostná
hodnota \(b\) 
je \(1\).}{Pravdivostná hodnota \(a\) 
je \(1\), pravdivostná
hodnota \(b\) 
je \(1\).}{Pravdivostná hodnota \(a\) 
je \(1\), pravdivostná
hodnota \(b\) 
je \(0\).}{Pravdivostná hodnota \(a\) 
je \(0\), pravdivostná
hodnota \(b\) 
je \(0\).}{}{}}
\LANGes{
\Question{
Determina los valores de verdad de las proposiciones  \(a\) y \(b\) sabiendo que la proposición compuesta
\[ 
 a \iff (a \vee b)
\]
es falsa.}
{La proposición  \(a\) 
es falsa, \(b\) 
es verdadera.}{Ambas proposiciones son verdaderas.}{La proposición \(a\) 
es verdadera, \(b\) 
es falsa.}{Ambas proposiciones son falsas.}{}{}}
\End

\Data\ID{9000086607}
\Updated{2020-07-15 03:07:37}
\Level{B}
\SubArea{Logic and Sets}
\LANGen{
\Question{
Determine the truth values of propositions \(a\) and \(b\) if you know that the compound proposition
\[ 
 (\neg a \vee b) \wedge a
\]
is true.}
{Both statements are true.}{The statement \(a\) 
is true, \(b\) 
is false.}{The statement \(a\) 
is false, \(b\) 
is true.}{Both statements are false.}{}{}}
\LANGpl{
\Question{
Oceń wartość logiczną zdań \(a\) 
i  \(b\), jeśli wiadomo, że zdanie
\[ 
 (\neg a \vee b) \wedge a
\]
jest prawdziwe.}
{Oba zdania są prawdziwe.}{Zdanie \(a\) jest prawdziwe, zdanie \(b\) jest fałszywe.}{Zdanie \(a\) jest fałszywe, zdanie \(b\) jest prawdziwe.}{Oba zdania są fałszywe.}{}{}}
\LANGcs{
\Question{
Určete pravdivostní hodnoty výroků
\(a\) a
\(b\),
víte-li, že pravdivostní hodnota složeného výroku
\((\neg a \vee b) \wedge a\) je
\(1\).}
{Pravdivostní hodnota \(a\) 
je \(1\), pravdivostní
hodnota \(b\) 
je \(1\).}{Pravdivostní hodnota \(a\) 
je \(1\), pravdivostní
hodnota \(b\) 
je \(0\).}{Pravdivostní hodnota \(a\) 
je \(0\), pravdivostní
hodnota \(b\) 
je \(1\).}{Pravdivostní hodnota \(a\) 
je \(0\), pravdivostní
hodnota \(b\) 
je \(0\).}{}{}}
\LANGsk{
\Question{
Určte pravdivostné hodnoty výrokov
\(a\) a
\(b\),
ak viete, že pravdivostná hodnota zloženého výroku
\((\neg a \vee b) \wedge a\) je
\(1\).}
{Pravdivostná hodnota \(a\) 
je \(1\), pravdivostná
hodnota \(b\) 
je \(1\).}{Pravdivostná hodnota \(a\) 
je \(1\), pravdivostná
hodnota \(b\) 
je \(0\).}{Pravdivostná hodnota \(a\) 
je \(0\), pravdivostná
hodnota \(b\) 
je \(1\).}{Pravdivostná hodnota \(a\) 
je \(0\), pravdivostná
hodnota \(b\) 
je \(0\).}{}{}}
\LANGes{
\Question{
Determina los valores de verdad de las proposiciones \(a\) y \(b\) sabiendo que la proposición compuesta
\[ 
 (\neg a \vee b) \wedge a
\]
es verdadera.}
{Ambas proposiciones son verdaderas.}{La proposición  \(a\) 
es verdadera, \(b\) 
es falsa.}{La proposición \(a\) 
es falsa, \(b\) 
es verdadera.}{Ambas proposiciones son falsas.}{}{}}
\End

\Data\ID{9000086608}
\Updated{2020-10-02 10:10:46}
\Level{B}
\SubArea{Logic and Sets}
\LANGen{
\Question{
Determine the truth values of statements \(a\) and \(b\) if you know that the compound statement
\[ 
 \neg a \iff (a \wedge b)
\]
is true.}
{The statement \(a\) 
is true, \(b\) 
is false.}{Both statements are true.}{The statement \(a\) 
is false, \(b\) 
is true.}{Both statements are false.}{}{}}
\LANGpl{
\Question{
Oceń wartość logiczną zdań  \(a\) 
i \(b\),
jeśli wiadomo, że zdanie
\[ 
 \neg a \iff (a \wedge b)
\]
jest prawdziwe.}
{Zdanie \(a\) jest prawdziwe, zdanie \(b\) jest fałszywe.}{Oba zdania są prawdziwe.}{Zdanie \(a\) jest fałszywe, zdanie \(b\) jest prawdziwe.}{Oba zdania są fałszywe.}{}{}}
\LANGcs{
\Question{
Určete pravdivostní hodnoty výroků
\(a\) a
\(b\),
víte-li, že pravdivostní hodnota složeného výroku
\(\neg a \iff (a \wedge b)\) je
\(1\).}
{Pravdivostní hodnota \(a\) 
je \(1\), pravdivostní
hodnota \(b\) 
je \(0\).}{Pravdivostní hodnota \(a\) 
je \(1\), pravdivostní
hodnota \(b\) 
je \(1\).}{Pravdivostní hodnota \(a\) 
je \(0\), pravdivostní
hodnota \(b\) 
je \(1\).}{Pravdivostní hodnota \(a\) 
je \(0\), pravdivostní
hodnota \(b\) 
je \(0\).}{}{}}
\LANGsk{
\Question{
Určte pravdivostné hodnoty výrokov
\(a\) a
\(b\),
ak viete, že pravdivostná hodnota zloženého výroku
\(\neg a \iff (a \wedge b)\) je
\(1\).}
{Pravdivostná hodnota \(a\) 
je \(1\), pravdivostná
hodnota \(b\) 
je \(0\).}{Pravdivostná hodnota \(a\) 
je \(1\), pravdivostná
hodnota \(b\) 
je \(1\).}{Pravdivostná hodnota \(a\) 
je \(0\), pravdivostná
hodnota \(b\) 
je \(1\).}{Pravdivostná hodnota \(a\) 
je \(0\), pravdivostná
hodnota \(b\) 
je \(0\).}{}{}}
\LANGes{
\Question{
Determina los valores de verdad de las proposiciones \(a\) y \(b\) sabiendo que la proposición compuesta 
\[ 
 \neg a \iff (a \wedge b)
\]
es verdadera.}
{La proposición \(a\) 
es verdadera, \(b\) 
es falsa.}{Ambas proposiciones son verdaderas.}{La proposición \(a\) 
es falsa, \(b\) 
es verdadera.}{Ambas proposiciones son falsas.}{}{}}
\End

\Data\ID{9000086609}
\Updated{2022-04-23 05:04:57}
\Level{B}
\SubArea{Logic and Sets}
\LANGen{
\Question{
Let \(a\) be a true statement while statements \(b\) and \(c\) are both false. Find the true statement.}
{\((a \vee b)\implies \neg c\)}{\((\neg a \vee b) \vee c\)}{\((a \wedge b) \vee c\)}{\(a \iff (b \vee c)\)}{}{}}
\LANGpl{
\Question{
Zdanie  \(a\) jest prawdziwe, natomiast zdania
 \(b\) 
i  \(c\) 
są fałszywe. Wskaż zdanie prawdziwe.}
{\((a \vee b)\implies \neg c\)}{\((\neg a \vee b) \vee c\)}{\((a \wedge b) \vee c\)}{\(a \iff (b \vee c)\)}{}{}}
\LANGcs{
\Question{
Je dán pravdivý výrok \(a\),
nepravdivý výrok \(b\) a
nepravdivý výrok \(c\).
Určete, který ze složených výroků je pravdivý.}
{\((a \vee b)\implies \neg c\)}{\((\neg a \vee b) \vee c\)}{\((a \wedge b) \vee c\)}{\(a \iff (b \vee c)\)}{}{}}
\LANGsk{
\Question{
Je daný pravdivý výrok \(a\),
nepravdivý výrok \(b\) a
nepravdivý výrok \(c\).
Určte, ktorý zo zložených výrokov je pravdivý.}
{\((a \vee b)\implies \neg c\)}{\((\neg a \vee b) \vee c\)}{\((a \wedge b) \vee c\)}{\(a \iff (b \vee c)\)}{}{}}
\LANGes{
\Question{
Sea \(a\) una proposición verdadera mientras que \(b\) y \(c\) son ambas falsas.}
{\((a \vee b)\implies \neg c\)}{\((\neg a \vee b) \vee c\)}{\((a \wedge b) \vee c\)}{\(a \iff (b \vee c)\)}{}{}}
\End

\Data\ID{9000086610}
\Updated{2020-10-02 10:10:14}
\Level{B}
\SubArea{Logic and Sets}
\LANGen{
\Question{
Let statements \(a\) 
and \(b\) be
false while \(c\) is a true statement. Find the false statement.}
{\(a \iff (b \vee c)\)}{\((\neg a \vee b) \vee c\)}{\((a \wedge b) \vee c\)}{\((a \vee b)\implies \neg c\)}{}{}}
\LANGpl{
\Question{
Zdania  \(a\) 
i  \(b\) są 
fałszywe,  natomiast zdanie \(c\) 
jest prawdziwe. Wskaż zdanie fałszywe.}
{\(a \iff (b \vee c)\)}{\((\neg a \vee b) \vee c\)}{\((a \wedge b) \vee c\)}{\((a \vee b)\implies \neg c\)}{}{}}
\LANGcs{
\Question{
Je dán nepravdivý výrok \(a\),
nepravdivý výrok \(b\) 
a pravdivý výrok \(c\).
Určete, který ze složených výroků je nepravdivý.}
{\(a \iff (b \vee c)\)}{\((\neg a \vee b) \vee c\)}{\((a \wedge b) \vee c\)}{\((a \vee b)\implies \neg c\)}{}{}}
\LANGsk{
\Question{
Je daný nepravdivý výrok \(a\),
nepravdivý výrok \(b\) 
a pravdivý výrok \(c\).
Určte, ktorý zo zložených výrokov je nepravdivý.}
{\(a \iff (b \vee c)\)}{\((\neg a \vee b) \vee c\)}{\((a \wedge b) \vee c\)}{\((a \vee b)\implies \neg c\)}{}{}}
\LANGes{
\Question{
Sean \(a\) 
y \(b\) proposiciones falsas mientras que \(c\) es una proposición verdadera.}
{\(a \iff (b \vee c)\)}{\((\neg a \vee b) \vee c\)}{\((a \wedge b) \vee c\)}{\((a \vee b)\implies \neg c\)}{}{}}
\End

\Data\ID{1003034401}
\Updated{2023-06-28 11:06:10}
\Level{C}
\SubArea{Logic and Sets}
\LANGen{
\Question{
There are \( 47 \) students in a class. Among these students \( 22 \) are members of a sports club, \( 33 \) are members of a film club and \( 20 \) students joined a literary club. Further, \( 17 \) students are members of the sports club and the film club, \( 13 \) joined the film club and the literary club, \( 6 \) attend the sports club and the literary club.  Only one student is a member of all the clubs.  How many students from the class attend neither the sports club nor the literary club?}
{\( 11 \)}{\( 7 \)}{\( 36 \)}{\( 4 \)}{}{}}
\LANGpl{
\Question{
Wśród \( 47 \) uczniów klas pierwszych pewnego liceum  \( 22 \)  uczniów należy do koła sportowego,  \( 33 \) do koła filmowego,  \( 20 \) uczniów do koła literackiego.Wśród tych uczniów \( 17 \) należy do koła sportowego i filmowego,  \( 13 \) 
 do filmowego i literackiego, \( 6 \) do sportowego i literackiego, tylko jeden uczeń należy do wszystkich kół. Oblicz ilu uczniów nie należy ani do koła sportowego ani do koła literackiego?}
{\( 11 \)}{\( 7 \)}{\( 36 \)}{\( 4 \)}{}{}}
\LANGcs{
\Question{
Ve třídě je  \( 47 \) studentů.  \( 22 \) z nich jsou členy sportovního klubu, \( 33 \)  jsou členy filmového klubu a \( 20 \) studentů se připojilo k literárnímu klubu. Dále, \( 17 \) studentů jsou členy sportovního i filmového klubu, \( 13 \) se připojilo k filmovému a literárnímu klubu, \( 6 \) navštěvuje sportovní i literární klub.  Jen jeden student je členem všech klubů. Kolik studentů ze třídy nenavštěvuje ani sportovní klub a ani literární klub?}
{\( 11 \)}{\( 7 \)}{\( 36 \)}{\( 4 \)}{}{}}
\LANGsk{
\Question{
V triede je  \( 47 \) študentov.  \( 22 \) z nich sú členmi športového klubu, \( 33 \)  sú členmi filmového klubu a \( 20 \) študentov sa zapojilo do literárneho klubu. Ďalej, \( 17 \) študentov je členom športového a aj filmového klubu, \( 13 \) sa zapojilo do filmového a literárneho klubu, \( 6 \) navštevujú športový a literárny klub. Iba jeden študent je členom všetkých klubov. Koľko študentov z triedy nenavštevuje ani športový klub a ani literárny klub?}
{\( 11 \)}{\( 7 \)}{\( 36 \)}{\( 4 \)}{}{}}
\LANGes{
\Question{
En una clase hay  \( 47 \) estudiantes. Entre estos estudiantes  \( 22 \) son miembros de un club deportivo, \( 33 \) son miembros de un club de cine y  \( 20 \) estudiantes se unieron a un club literario. Además, \( 17 \) estudiantes son miembros del club deportivo y del club de cine, \( 13 \) se unieron al club de cine y al club literario,  \( 6 \) acuden al club deportivo y al club literario. Solo un estudiante es miembro de todos los clubs. ¿Cuántos estudiantes de la clase no acuden al club deportivo ni al club literario?}
{\( 11 \)}{\( 7 \)}{\( 36 \)}{\( 4 \)}{}{}}
\End

\Data\ID{1003034402}
\Updated{2022-08-25 10:08:03}
\Level{C}
\SubArea{Logic and Sets}
\LANGen{
\Question{
In a  class of \( 31 \) pupils, \( 18 \) pupils learn Italian, \( 16 \) pupils learn Spanish and \( 12 \) pupils learn Portuguese. Now, \( 7 \) pupils learn Italian and Spanish, \( 9 \) pupils learn  Italian and Portuguese and \( 5 \) pupils learn Spanish and Portuguese. Finally, \( 3 \) students do not learn any of these three languages. How many pupils learn all these languages?}
{\( 3 \)}{\( 5 \)}{\( 2 \)}{\( 1 \)}{}{}}
\LANGpl{
\Question{
W klasie liczącej  \( 31 \) osób, \( 18 \) uczniów uczy się języka włoskiego, \( 16 \) uczniów uczy się języka hiszpańskiego i \( 12 \) uczniów uczy się języka portugalskiego. Wśród nich  \( 7 \) uczniów uczy się dwóch języków włoskiego i hiszpańskiego, \( 9 \)  uczniów uczy się języka włoskiego i portugalskiego i  \( 5 \) uczniów uczy się języka hiszpańskiego i portugalskiego. Wreszcie  \( 3 \)  uczniów nie uczy się żadnego z tych trzech języków. Ilu uczniów uczy się wszystkich języków?}
{\( 3 \)}{\( 5 \)}{\( 2 \)}{\( 1 \)}{}{}}
\LANGcs{
\Question{
Z \( 31 \) žáků třídy se \( 18 \) učí italštinu, \( 16 \) španělštinu a \( 12 \) žáků se učí portugalštinu. Dále víme, že \( 7 \) žáků se učí italštinu a španělštinu, \( 9 \) žáků se učí italštinu a portugalštinu a \( 5 \) žáků se učí španělštinu a portugalštinu. Nakonec, \( 3 \) studenti se neučí žádný z těchto třech jazyků. Určete, kolik žáků se učí všechny tyto jazyky.}
{\( 3 \)}{\( 5 \)}{\( 2 \)}{\( 1 \)}{}{}}
\LANGsk{
\Question{
Z \( 31 \) žiakov triedy, \( 18 \) sa učí taliančinu, \( 16 \) žiakov sa učí španielčinu a \( 12 \) žiakov sa učí portugalčinu.Teraz \( 7 \) žiakov sa učí taliančinu a španielčinu, \( 9 \) žiakov sa učí taliančinu a portugalčinu a \( 5 \) žiakov sa učí španielčinu a portugalčinu. Nakoniec, \( 3 \) študenti sa neučia žiaden z týchto troch jazykov. Koľko žiakov sa učí všetky tieto jazyky?}
{\( 3 \)}{\( 5 \)}{\( 2 \)}{\( 1 \)}{}{}}
\LANGes{
\Question{
En una clase de  \( 31 \) alumnos, \( 18 \) alumnos aprenden italiano, \( 16 \) alumnos aprenden español y  \( 12 \) alumnos aprenden portugués. Ahora, \( 7 \) alumnos aprenden italiano y español,  \( 9 \) alumnos aprenden italiano y portugués y  \( 5 \) alumnos aprenden español y portugués. Finalmente, \( 3 \) estudiantes no aprenden ninguno de estos tres idiomas. ¿Cuántos alumnos aprenden todos estos idiomas?}
{\( 3 \)}{\( 5 \)}{\( 2 \)}{\( 1 \)}{}{}}
\End

\Data\ID{1003034403}
\Updated{2023-06-29 01:06:03}
\Level{C}
\SubArea{Logic and Sets}
\LANGen{
\Question{
In a  group of \( 100 \) people, \( 60 \) people like swimming, \( 50 \) people like horse-riding, \( 36 \) people like cycling. Further, \( 16 \) people like swimming and horse riding, \( 19 \) people like swimming and cycling, \( 15 \) people like horse-riding and cycling.  How many people in the group enjoy only one of these activities?}
{\( 58 \)}{\( 29 \)}{\( 23 \)}{\( 6 \)}{}{}}
\LANGpl{
\Question{
W grupie  \( 100 \) osób, \( 60 \) lubi pływać, \( 50 \) lubi jazdę konną, \( 36 \)  jeździ na rowerze, \( 16 \) lubi pływać i jazdę konną, \( 19 \) lubi pływać i jazdę na rowerze,  \( 15 \) jazdę konną i rower. Ile osób lubi tylko jedną z tych trzech aktywności?}
{\( 58 \)}{\( 29 \)}{\( 23 \)}{\( 6 \)}{}{}}
\LANGcs{
\Question{
Ve skupině \( 100 \) lidí, \( 60 \) lidí rádo plave, \( 50 \) lidí rádo jezdí na koni, \( 36 \) lidí rádo jezdí na kole. Dále, \( 16 \) lidí rádo plave a i rádo jezdí na koni, \( 19 \) lidí rádo plave a i jezdí na kole, \( 15 \) lidí rádo jezdí na koni a na kole. Kolik lidí z této skupiny má rádo jen jednu z těchto aktivit?}
{\( 58 \)}{\( 29 \)}{\( 23 \)}{\( 6 \)}{}{}}
\LANGsk{
\Question{
V skupine \( 100 \) ľudí, \( 60 \) ľudí pláva, \( 50 \) ľudí jazdí na koni, \( 36 \) ľudí bicykluje. Ďalej, \( 16 \) ľudí pláva a aj jazdí na koni, \( 19 \) ľudí pláva a aj jazdí na bicykli, \( 15 \) ľudí jazdí na koni a bicykluje. Koľko ľudí z tejto skupiny realizuje len jednu z týchto aktivít?}
{\( 58 \)}{\( 29 \)}{\( 23 \)}{\( 6 \)}{}{}}
\LANGes{
\Question{
En un grupo de \( 100 \) personas, a  \( 60 \) personas les gusta nadar, a \( 50 \) personas les gusta montar a caballo y a \( 36 \) personas les gusta montar en bicicleta. Además, a  \( 16 \) personas les gusta nadar y montar a caballo, a  \( 19 \) personas les gusta nadar y montar en bicicleta y a \( 15 \) personas les gusta montar a caballo y montar en bicicleta. ¿Cuántas personas en el grupo disfrutan practicando solo una de estas actividades?}
{\( 58 \)}{\( 29 \)}{\( 23 \)}{\( 6 \)}{}{}}
\End

\Data\ID{2010001406}
\Updated{2023-03-20 04:03:41}
\Level{C}
\SubArea{Logic and Sets}
\LANGen{
\Question{
All \(96\) children in the summer camp were to choose the chocolate champion - between milk chocolate and dark chocolate. At the end milk chocolate got \(30\) more votes than dark chocolate. \(26\) children couldn’t decide though and gave their votes to both chocolates. How many votes did dark chocolate get?}
{\(46\)}{\( 76\)}{\( 26\)}{\( 82\)}{}{}}
\LANGpl{
\Question{
Wszystkie \(96\) dzieci na letnim obozie miało wybrać czekoladowego mistrza - między mleczną a gorzką czekoladą. Ostatecznie czekolada mleczna zdobyła \(30\) głosów więcej niż czekolada gorzka. \(26\)  dzieci nie mogło się jednak zdecydować i oddało głosy na obie czekolady. Ile głosów otrzymała ciemna czekolada?}
{\(46\)}{\( 76\)}{\( 26\)}{\( 82\)}{}{}}
\LANGcs{
\Question{
Všech \(96\) dětí na letním táboře se zúčastnilo ankety, jaký druh čokolády mají nejraději, zda hořkou nebo mléčnou. Mléčná čokoláda získala o \(30\) hlasů více než hořká čokoláda. \(26\) dětí se však nedokázalo rozhodnout a dalo hlas oběma čokoládám.. Kolik hlasů získala hořká čokoláda?}
{\(46\)}{\( 76\)}{\( 26\)}{\( 82\)}{}{}}
\LANGsk{
\Question{
Všetkých \(96\) detí v letnom tábore sa zúčastnilo ankety, aký druh čokolády majú najradšej, či horkú alebo mliečnu. Mliečna čokoláda získala o 30
 hlasov viac ako horká čokoláda. \(26\) detí sa však nevedelo rozhodnúť a dali hlas obidvom čokoládam. Koľko hlasov získala horká čokoláda?}
{\(46\)}{\( 76\)}{\( 26\)}{\( 82\)}{}{}}
\LANGes{
\Question{
Los \(96\) niños de un campamento de verano participaron en la encuesta sobre su chocolate favorito. Las opciones eran chocolate con leche y chocolate negro. Al final, el chocolate con leche obtuvo \(30\) votos más que el chocolate negro. \(26\) niños no pudieron decidir y dieron su voto a ambos chocolates. ¿Cuántos votos obtuvo el chocolate negro?}
{\(46\)}{\( 76\)}{\( 26\)}{\( 82\)}{}{}}
\End

\Data\ID{9000089001}
\Updated{2022-04-23 05:04:57}
\Level{C}
\SubArea{Logic and Sets}
\LANGen{
\Question{
Students from a class have a possibility to work in mathematical and physical hobby groups. There are \(31\) students in the class. From the total, \(21\) students are members of mathematical group. Some of the students are members of both groups, but there are \(10\) which are members of just one group. There are \(3\) students which are not members of any of those groups. How many students are members of both mathematical and physical groups?}
{\(18\)}{\(16\)}{\(19\)}{}{}{}}
\LANGpl{
\Question{
Grupa 31 uczniów ma możliwość uczęszczania na kółko matematyczne i fizyczne. \(21\) uczęszcza na kółko matematyczne. \(10\) uczęszcza na jedno kółko. \(3\) nie uczęszcza nigdzie. Ilu uczniów uczęszcza na oba kółka?}
{\(18\)}{\(16\)}{\(19\)}{}{}{}}
\LANGcs{
\Question{
Žáci 1. ročníku mohou navštěvovat matematický a fyzikální kroužek. Z \(31\) žáků třídy jich \(21\) chodí do matematického kroužku. Pouze do jednoho z kroužků chodí \(10\) žáků. Do žádného kroužku nechodí \(3\) žáci. Kolik žáků chodí do obou kroužků zároveň?}
{\(18\)}{\(16\)}{\(19\)}{}{}{}}
\LANGsk{
\Question{
Žiaci 1. ročníka môžu navštevovať matematický a fyzikálny krúžok. Z
\(31\) žiakov
triedy  \(21\) 
chodí do matematického krúžku. Len do jedného z krúžku chodí
\(10\) žiakov. Do žiadneho
krúžku nechodia \(3\) 
žiaci. Koľko žiakov chodí do oboch krúžkov zároveň?}
{\(18\)}{\(16\)}{\(19\)}{}{}{}}
\LANGes{
\Question{
Los estudiantes de una clase tienen la posibilidad de trabajar en grupos matemáticos y físicos según sus aficiones. Hay  \(31\) estudiantes en la clase. Del total,  \(21\) estudiantes son miembros del grupo matemático. Algunos de los estudiantes son miembros de ambos grupos, pero hay  \(10\) que son miembros de un solo grupo. Hay \(3\) estudiantes que no son miembros de ninguno de esos grupos. ¿Cuántos estudiantes son miembros de ambos grupos (matemático y físico)?}
{\(18\)}{\(16\)}{\(19\)}{}{}{}}
\End

\Data\ID{9000089002}
\Updated{2023-06-28 11:06:58}
\Level{C}
\SubArea{Logic and Sets}
\LANGen{
\Question{
Students from a class decided to order books for forthcoming holiday. The book-shop in the neighborhood had two bestsellers on the stock: a crime novel and a horror stories. There were \(31\) students in the class in total. From this total, \(22\) students bought horror stories. Altogether \(12\) students bought just one of the books. Two students did not buy any of these books. How many students bought the crime novel?}
{\(24\)}{\(7\)}{\(5\)}{}{}{}}
\LANGpl{
\Question{
Uczniowie pewnej klasy postanowili zakupić książki na nadchodzące wakacje w pobliskiej księgarni, w której są dwa rodzaje książek: powieść kryminalna i horror. Jest \(31\) uczniów w klasie. \(22\) uczniów kupiło horror. \(12\) uczniów kupiło jedną z tych książek. Dwóch uczniów nie kupiło żadnej książki. Ilu uczniów kupiło powieść kryminalną?}
{\(24\)}{\(7\)}{\(5\)}{}{}{}}
\LANGcs{
\Question{
Žáci 1. ročníku si kupovali knížky, aby se nenudili o
nadcházejících prázdninách. V blízké prodejně právě
dostali dlouho očekávanou detektivku a strašidelný horor. Z
\(31\) žáků
třídy si \(22\) 
koupilo horor. Pouze jednu z těchto knih si koupilo
\(12\) 
žáků. Žádnou z těchto knih si nekoupili dva žáci. Kolik žáků si koupilo
detektivku?}
{\(24\)}{\(7\)}{\(5\)}{}{}{}}
\LANGsk{
\Question{
Žiaci 1. ročníka si kupovali knižky, aby sa nenudili cez
blížiace sa prázdniny. V blízkej predajni práve
dostali dlho očakávanú detektívku a strašidelný horor. Z
\(31\) žiakov
triedy si \(22\) 
kúpilo horor. Iba jednu z týchto kníh si kúpilo
\(12\) 
žiakov. Žiadnu z týchto kníh si nekúpili dvaja žiaci. Koľko žiakov si kúpilo
detektívku?}
{\(24\)}{\(7\)}{\(5\)}{}{}{}}
\LANGes{
\Question{
Los estudiantes de una clase decidieron comprar libros para las próximas vacaciones. La librería del barrio tenía dos superventas en el almacén: una novela policíaca y una historia de terror. En total había \(31\) estudiantes en la clase. De este total, \(22\) estudiantes compraron la historia de terror. En total  \(12\) estudiantes compraron solo un libro. Dos estudiantes no compraron ningún libro. ¿Cuántos estudiantes compraron la novela policíaca?}
{\(24\)}{\(7\)}{\(5\)}{}{}{}}
\End

\Data\ID{9000089003}
\Updated{2023-06-29 01:06:50}
\Level{C}
\SubArea{Logic and Sets}
\LANGen{
\Question{
Students from a class bought a snack in the school lunch-room. There were \(31\) students in the class in total. Altogether \(8\) students had snack from their home and they did not buy anything. Altogether \(12\) students bought hamburger and \(15\) students bought hot-dog. How many students bought both hamburger and hot-dog?}
{\(4\)}{\(19\)}{\(8\)}{}{}{}}
\LANGpl{
\Question{
Uczniowie pierwszej klasy mogli sobie kupić przekąskę w szkolnej stołówce. Jest \(31\) uczniów w klasie. \(8\) uczniów przyniosło przekąskę z domu, więc nic nie kupili. \(12\) uczniów kupiło hamburgera, natomiast \(15\) uczniów kupiło hot-doga. Ilu uczniów kupiło zarówno hamburgera jak i hot-doga?}
{\(4\)}{\(19\)}{\(8\)}{}{}{}}
\LANGcs{
\Question{
Žáci 1. ročníku nakupovali svačinku ve školním bufetu. Z
\(31\) žáků mělo
\(8\) svačinku z domu, a proto
si nic nekoupilo. \(12\) dětí si
koupilo housku se sekanou a \(15\) 
žáků si koupilo housku s vuřtem. Kolik nenasytů si koupilo oba typy housky?}
{\(4\)}{\(19\)}{\(8\)}{}{}{}}
\LANGsk{
\Question{
Žiaci 1. ročníka nakupovali desiatu v školskom bufete. Z
\(31\) žiakov malo
\(8\) desiatu z domu, a preto
si nič nekúpili. \(12\) detí si
kúpilo bagetu so šunkou a \(15\) 
žiakov si kúpilo bagetu so zeleninou. Koľko hladošov si kúpilo oba typy bagiet?}
{\(4\)}{\(19\)}{\(8\)}{}{}{}}
\LANGes{
\Question{
Los estudiantes de una clase compraron el almuerzo en el comedor escolar. Había un total de \(31\) estudiantes en la clase. De ese total,  \(8\) estudiantes se tomaron el almuerzo traido de su casa y no compraron nada. En total, \(12\) estudiantes compraron hamburguesas y \(15\) estudiantes compraron perritos calientes. ¿Cuántos estudiantes compraron hamburguesas junto con perritos calientes?}
{\(4\)}{\(19\)}{\(8\)}{}{}{}}
\End

\Data\ID{9000089004}
\Updated{2023-06-28 11:06:55}
\Level{C}
\SubArea{Logic and Sets}
\LANGen{
\Question{
Students in a university are allowed to buy either lunch or dinner in a student's canteen. There are \(129\) freshmen students in total. Altogether \(116\) freshmen students buy the lunch or dinner. From this amount, \(62\) students buy just one meal. The number of freshmen students which buy lunch is bigger by
 \(46\) than the number of freshmen students which buy dinner. How many freshmen students buy only the dinner?}
{\(8\)}{\(54\)}{\(62\)}{}{}{}}
\LANGpl{
\Question{
Uczniowie pierwszego roku mogą sobie kupić  drugie śniadanie i  obiad na stołówce. Jest \(129\) studentów pierwszego roku. \(116\) studentów kupuje drugie śniadanie lub obiad.  \(62\) studentów kupuje tylko jeden z tych posiłków. Ponadto  drugie śniadanie kupuje o \(46\) studentów więcej niż obiad. Ilu studentów pierwszego roku kupuje  tylko na obiad?}
{\(8\)}{\(54\)}{\(62\)}{}{}{}}
\LANGcs{
\Question{
Ze \(129\) 
studentů prvního ročníku vysoké školy chodí do menzy na oběd nebo
večeři \(116\) 
studentů, \(62\) 
studentů dochází právě na jedno z těchto jídel. Přitom na obědy chodí
o \(46\) 
studentů více než na večeře. Kolik studentů prvního ročníku chodí
jenom na večeře?}
{\(8\)}{\(54\)}{\(62\)}{}{}{}}
\LANGsk{
\Question{
Zo \(129\) 
študentov prvého ročníka vysokej školy chodí do menzy na obed alebo
večeru \(116\) 
študentov, \(62\) 
študentov prichádza práve na jedno z týchto jedál. Pritom na obedy chodí
o \(46\) 
študentov viac než na večere. Koľko študentov prvého ročníka chodí
len na večere?}
{\(8\)}{\(54\)}{\(62\)}{}{}{}}
\LANGes{
\Question{
Los estudiantes de una universidad pueden comprar el almuerzo o la cena en el comedor escolar. En total hay \(129\) estudiantes del primer curso. Un  total de \(116\) estudiantes del primer curso compran la comida o la cena. De esta cantidad,  \(62\) estudiantes compran solo una comida. El número de estudiantes del primer curso que compran almuerzo es mayor en 
 \(46\) unidades al número de estudiantes del primer curso que compran la cena. ¿Cuántos estudiantes del primer curso compran solo la cena?}
{\(8\)}{\(54\)}{\(62\)}{}{}{}}
\End

\Data\ID{9000089005}
\Updated{2023-06-28 11:06:02}
\Level{C}
\SubArea{Logic and Sets}
\LANGen{
\Question{
There are two types of cheese in a shop. Altogether \(153\) customer were in the shop during a day. From this total, \(65\) customers bought the first cheese. From the same total, \(49\) customers bought the second cheese. \(20\%\) of customers which bought at least one of the types of cheese bought actually both types. How many customers did not buy any of these two cheeses?}
{\(58\)}{\(39\)}{\(19\)}{}{}{}}
\LANGpl{
\Question{
W sklepie są dwa rodzaje sera. \(153\) klientów odwiedza sklep w ciągu dnia. \(65\) klientów  kupiło pierwszy rodzaj sera. Drugi rodzaj sera zakupiło \(49\) klientów. \(20\%\) klientów, którzy kupili jeden rodzaj sera kupiło oba rodzaje sera. Ilu klientów nie kupiło sera?}
{\(58\)}{\(39\)}{\(19\)}{}{}{}}
\LANGcs{
\Question{
V obchodě se objevily dva nové druhy sýrů. Ze
\(153\) zákazníků
jich \(65\) 
neodolalo koupi prvního druhu. Druhý druh zakoupilo
\(49\) 
zákazníků. Těch, kteří zakoupili oba druhy, byla pouze pětina počtu těch
zákazníků, kteří zakoupili aspoň jeden druh. Kolik zákazníků si
nekoupilo žádný z těchto sýrů?}
{\(58\)}{\(39\)}{\(19\)}{}{}{}}
\LANGsk{
\Question{
V obchode sa objavili dva nové druhy syra. Zo
\(153\) zákazníkov
ich \(65\) 
neodolalo kúpiť si prvý druh. Druhý druh si kúpilo
\(49\) 
zákazníkov. Tých, ktorí si kúpili oba druhy, bola len pätina počtu tých
zákazníkov, ktorí si kúpili aspoň jeden druh. Koľko zákazníkov si
nekúpilo žiadny z týchto syrov?}
{\(58\)}{\(39\)}{\(19\)}{}{}{}}
\LANGes{
\Question{
Hay dos tipos de queso en una tienda. En total,  \(153\) clientes visitaron la tienda durante un día. De ese total, \(65\) clientes compraron el primer queso. Del mismo total,  \(49\) clientes compraron el segundo queso. El  \(20\%\) de los clientes que compraron por lo menos uno de los tipos de queso compraron en realidad ambos tipos. ¿Cuántos clientes no compraron ninguno de estos dos quesos?}
{\(58\)}{\(39\)}{\(19\)}{}{}{}}
\End

\Data\ID{9000089006}
\Updated{2023-06-28 11:06:14}
\Level{C}
\SubArea{Logic and Sets}
\LANGen{
\Question{
A questionnaire on popular singers has been filled by \(200\) girls from a high-school. The girls had to mark the singers from three most popular (Justin Timberlake, Justin Bieber and Axl Rose) which they like. Justin Timberlake got \(78\) votes, Justin Bieber got \(75\) votes and Axl Rose \(101\) votes. There were \(28\) girls which voted for all three singers. There were \(22\) girls which voted for two singers. One half of this amount are fans of the pair Bieber and Rose. The number of the girls which like only Justin Bieber is smaller by \(7\) comparing to the number of girls which like only Justin Timberlake. How many girls did not like any of these three singers?}
{\(24\)}{\(32\)}{\(11\)}{}{}{}}
\LANGpl{
\Question{
\(200\)  uczennic liceum wypełniło ankietę dotyczącą ulubionych piosenkarzy (Justin Timberlake, Justin Bieber and Axl Rose). Justin Timberlake otrzymał \(78\) głosów, Justin Bieber otrzymał \(75\) głosów, natomiast Axl Rose otrzymał \(101\) głosów. \(28\) dziewczyn zagłosowało na każdego z nich. \(22\) 
 dziewczyny zagłosowały na dwóch piosenkarzy. Połowa z nich to fanki pary Bieber i Rose. Dziewczyn, które lubią tylko Justina Biebera jest o \(7\) mniej niż tych, które lubią tylko Justina Timberlakea. Ile dziewczyn nie lubi żadnego z tych piosenkarzy?}
{\(24\)}{\(32\)}{\(11\)}{}{}{}}
\LANGcs{
\Question{
Celkem \(200\) 
studentek gymnázia vyplnilo anketu s otázkou, který ze tří
zpěváků (K. Gott, M. David, D. Hůlka) se jim líbí. Gotta obdivuje
\(78\), Davida
\(75\) a Hůlku
\(101\) 
děvčat. Všechny tři zpěváky současně obdivuje
\(28\) 
studentek. Těch, které obdivují právě dva z těchto zpěváků, je
\(22\) a z nich
právě polovinu tvoří obdivovatelky dvojice David, Hůlka. Děvčat, která obdivují
jen Davida, je o \(7\) 
méně než těch, která obdivují jen Gotta. Kolik děvčat neobdivuje nikoho
z těchto tří populárních zpěváků?}
{\(24\)}{\(32\)}{\(11\)}{}{}{}}
\LANGsk{
\Question{
Celkom \(200\) študentiek gymnázia vyplnilo anketu s otázkou, ktorý z troch spevákov (K. Gott, M. David, D. Hůlka) sa im páči. Gotta obdivuje \(78\), Davida \(75\) a Hůlku \(101\) dievčat. Všetkých troch spevákov súčasne obdivuje \(28\) študentiek. Tých, ktoré obdivujú práve dvoch z týchto spevákov, je \(22\) a z nich
práve polovicu tvoria obdivovateľky dvojice David, Hůlka. Dievčatá, ktoré obdivujú len Davida, je o \(7\) menej než tých, ktoré obdivujú jen Gotta. Koľko dievčat neobdivuje nikoho z týchto troch populárnych spevákov?}
{\(24\)}{\(32\)}{\(11\)}{}{}{}}
\LANGes{
\Question{
Un cuestionario sobre cantantes populares ha sido completado por \(200\) chicas de una escuela secundaria. Las chicas tuvieron que votar a tres de los cantantes más populares (Justin Timberlake, Justin Bieber y Axl Rose) que les gustan. Justin Timberlake obtuvo  \(78\) votos, Justin Bieber obtuvo \(75\) votos y Axl Rose  \(101\) votos. Había  \(28\) chicas que votaron por los tres cantantes y \(22\) chicas que votaron por dos cantantes. La mitad de esta cantidad son aficionadas a la pareja Bieber y Rose. El número de chicas a las que solo les gusta Justin Bieber es menor en  \(7\) unidades si lo comparamos con el número de chicas a las que solo les gusta Justin Timberlake. ¿A cuántas chicas no les gustó ninguno de estos tres cantantes?}
{\(24\)}{\(32\)}{\(11\)}{}{}{}}
\End

\Data\ID{9000089007}
\Updated{2020-07-15 03:07:37}
\Level{C}
\SubArea{Logic and Sets}
\LANGen{
\Question{
A class in the school has \(35\) students. On the last holiday the students visited Slovakia, Croatia and Bulgaria. From the total amount of \(35\), \(7\) students have been in Slovakia, \(7\) students have been in Croatia, \(5\) students have been in Bulgaria, \(21\) students have not been abroad, one student was in every of these countries, two students have been in both Bulgaria and Croatia but not in Slovakia, one student has been in both Bulgaria and Slovakia but not in Croatia. How many students visited either Slovakia or Croatia?}
{\(11\)}{\(7\)}{\(3\)}{}{}{}}
\LANGpl{
\Question{
Klasa liczy  \(35\) 
uczniów.  W wakacje uczniowie zwiedzali Słowację, Chorwację i Bułgarię. Z  \(35\) uczniów \(7\) uczniów było na wycieczce na Słowacji, \(7\) uczniów pojechało na Chorwację, \(5\) uczniów odwiedziło Bułgarię, \(21\) uczniów nigdzie nie wyjechało, jeden uczeń odwiedził wszystkie trzy kraje, dwóch uczniów było na Chorwacji i w Bułgarii, jeden uczeń był w Bułgarii i na Słowacji. Ilu uczniów zwiedziło Słowację lub Chorwację?}
{\(11\)}{\(7\)}{\(3\)}{}{}{}}
\LANGcs{
\Question{
Z \(35\) žáků 1.A bylo
\(7\) o prázdninách
na Slovensku, \(7\) v
Chorvatsku a \(5\) v
Bulharsku. Celkem \(21\) 
žáků o prázdninách do ciziny vůbec nevyjelo. Všechny tři krajiny
navštívil jeden žák. V Chorvatsku i v Bulharsku byli dva žáci. V Bulharsku i
na Slovensku byl jeden žák. Kolik žáků navštívilo přes prázdniny
Slovensko nebo Chorvatsko?}
{\(11\)}{\(7\)}{\(3\)}{}{}{}}
\LANGsk{
\Question{
Z \(35\) žiakov 1.A bolo
\(7\) cez prázdniny
vo Švajčiarsku, \(7\) v
Chorvátsku a \(5\) v
Bulharsku. Celkom \(21\) 
žiakov cez prázdniny do zahraničia vôbec nešlo. Všetky tri krajiny
navštívil jeden žiak. V Chorvátsku i v Bulharsku boli dvaja žiaci. V Bulharsku i
vo Švajčiarsku bol jeden žiak. Koľko žiakov navštívilo cez prázdniny
Švajčiarsko alebo Chorvátsko?}
{\(11\)}{\(7\)}{\(3\)}{}{}{}}
\LANGes{
\Question{
Una clase de la escuela tiene  \(35\) estudiantes. Durante las últimas vacaciones, los estudiantes visitaron Eslovaquia, Croacia y Bulgaria. De los  \(35\) estudiantes, \(7\) estudiantes estuvieron en Eslovaquia, \(7\) estudiantes en Croacia, \(5\) estudiantes en Bulgaria, \(21\) estudiantes no fueron al extranjero, un estudiante estuvo en todos estos países, dos estudiantes estuvieron en Bulgaria y en Croacia pero no en Eslovaquia, un estudiante estuvo en Bulgaria y en Eslovaquia pero no en Croacia. ¿Cuántos estudiantes visitaron Eslovaquia o Croacia?}
{\(11\)}{\(7\)}{\(3\)}{}{}{}}
\End
